% Experiments and results.
% Fragment: % Experiments and results.
% Fragment: % Experiments and results.
% Fragment: % Experiments and results.
% Fragment: \input{experiments} from the main article. Requires \usepackage{booktabs}.
% Every \tbd marks a result that has not been computed yet.
% Sources of the figures: docs/article/outline.md and, under method2-deterministic/,
% docs/evaluation/*.md, eval/eval_remote_1/eval/summary.txt, eval/eval_remote_2/eval/summary.txt.
% Method 1 full run: method1-rag/data/output/vps_run_2026-11-19/ (plan, audit report, run_all.log).
\providecommand{\tbd}{\textit{tbd}}

\section{Experiments and Results}
\label{sec:experiments}

Three methods are compared on the same task, generating a conformance test
plan from the text of a standard: a retrieval-augmented agentic pipeline
(Section~\ref{sec:exp-rag}), a deterministic pipeline in which the language
model only fills the steps a fixed procedure assigns to it
(Section~\ref{sec:exp-deterministic}), and an industrial agentic harness
given the task directly (Section~\ref{sec:exp-harness}). Each subsection
gives the method in a short paragraph and then its results;
Section~\ref{sec:exp-comparison} places the three side by side.

\subsection{Experimental setting}
\label{sec:exp-setting}

\paragraph{Input and ground truth.}
All experiments use the same input, the 126-page draft ISO/IEC DIS
18013-5:2020, and the same ground truth, the three appendices of ISO/IEC TS
18013-6:2025, parsed by rule into structured test cases
(Table~\ref{tab:ground-truth}). The 687 official cases that cite the standard
refer to 54 distinct clauses of the published 2021 edition, whereas the input
is the 2020 draft; this edition mismatch is a confound in every coverage
figure below and is quantified in Section~\ref{sec:exp-conformance}.

\begin{table}[t]
  \centering
  \caption{Ground truth: official test cases of ISO/IEC TS 18013-6:2025.}
  \label{tab:ground-truth}
  \begin{tabular}{lrr}
    \toprule
    Appendix & Test cases & Citing an 18013-5 clause \\
    \midrule
    1 -- Common and certificates & 138 & 14 \\
    2 -- mDL                     & 460 & 451 \\
    3 -- mDL reader              & 309 & 222 \\
    \midrule
    Total                        & 907 & 687 \\
    \bottomrule
  \end{tabular}
\end{table}

\paragraph{Metrics.}
\emph{Judge coverage} is the share of official cases for which an LLM judge,
shown the five generated cases closest to the official one, finds at least one
that verifies the same check (same data item or behaviour, same condition,
same expected outcome). \emph{Structural match} uses no model: each official
case is reduced by rule to a pair (data item, check type), and it is matched
when the generated plan carries a coverage item whose source object is named
by that data item and whose check belongs to the same family; 747 of the 907
official cases can be reduced to such a pair and are called \emph{scorable}.
\emph{Precision} is judged on a fixed random sample of 60 generated cases
against three criteria: the case is \emph{grounded} in its requirement,
\emph{actionable} by a test engineer, and its expected result is
\emph{correct}. The judge belongs to the same model family as the generator,
which is a known self-preference risk; the structural matcher was introduced
to remove it from the headline figure.

% ---------------------------------------------------------------------------
\subsection{Method 1: RAG pipeline}
\label{sec:exp-rag}

\paragraph{Method.}
The first method is an agentic retrieval-augmented pipeline organised in
three stages. Axis~A extracts the PDF, splits it with a structure-aware
two-pass chunker and indexes the chunks for hybrid dense and sparse retrieval
fused by reciprocal rank. Axis~B runs a Planner and then a chain of three
agents (Analyst, RAG Router, Designer) once per normative chunk, with
schema-constrained generation, an error-reflective retry and a BERTScore
self-evaluation gate, and assembles the output as an ISO/IEC/IEEE 29119-3
test plan. Axis~C audits each generated case with structural guardrails, an
embedding-distance pre-audit and an Auditor agent whose verdict is overridden
by a deterministic rule. The pipeline was developed with \texttt{llama3.2};
the full run reported here uses \texttt{gemma4:31b} for every agent, served
locally with a context of 4096 tokens.

\paragraph{Results: intermediate artefacts.}
The first documented runs produced chunks and an index but no test plan
(Table~\ref{tab:agentic-baseline}). The measurements of
Table~\ref{tab:diagnosis} show that a single defect, a heading detector that
matched object identifiers in the annexes instead of clause numbers, silently
disabled the chunking, cross-reference, defined-term and planning mechanisms.
A comparison of extraction back-ends (Table~\ref{tab:backends}) located the
cause in the extraction layer.

\paragraph{Results: the full run.}
One full run of the three axes was then made with the heading defect still
present (Table~\ref{tab:agentic-run}). Axis~B wrote a test case for 337 of
the 472 normative chunks (71.4\,\%); the other 135 were dropped because an
agent did not return valid JSON within its retries, the RAG Router alone
accounting for 102 of them. The loss does not depend on the strength of the
requirement (72\,\% of the chunks with a \emph{shall}, 69\,\% with a
\emph{should}, 71\,\% with a \emph{may}) but on the place in the document: it
is highest on the certificate-profile tables. Axis~C passed 154 of the 337
cases (45.7\,\%). Its hallucination and omission rates miss their targets and
its contradiction rate meets it (Table~\ref{tab:agentic-pending}). These rates
are an optimistic estimate: the Auditor is the model that wrote the cases, and
the six cases it failed to audit are counted as free of findings. The schema
check found no violation, and the embedding-distance pre-audit flagged no
case, its score being the same for passed and failed cases (0.84).

The heading defect is visible in the plan. The 1,246 chunks carry 27 distinct
section labels, so the 337 cases share 97 identifiers, 127 of them a single
one, the five feature sets have no meaningful name, and a case can only be
traced to its chunk, not to a clause. For that reason coverage is reported
per normative chunk and not per section, and the run has not been scored
against the official suite: both the judge and the structural matcher need a
clause or a named object on the generated side. The cases are also short,
261 of 337 having exactly two steps.

\begin{table}[t]
  \centering
  \caption{RAG pipeline: full run on ISO/IEC DIS 18013-5
    (\texttt{gemma4:31b}).}
  \label{tab:agentic-run}
  \begin{tabular}{lr}
    \toprule
    Measure & Value \\
    \midrule
    Chunks indexed (Axis A)                 & 1,246 \\
    \quad normative                         & 472 \\
    Distinct section labels                 & 27 \\
    \midrule
    Test cases generated (Axis B)           & 337 / 472 = 71.4\,\% \\
    \quad \emph{shall} / \emph{should} / \emph{may} & 248 / 61 / 28 \\
    Chunks dropped                          & 135 \\
    \quad RAG Router / Designer / Analyst   & 102 / 31 / 2 \\
    Distinct test case identifiers          & 97 \\
    Steps per case, mean                    & 2.3 \\
    \midrule
    Audit verdict (Axis C): pass            & 154 = 45.7\,\% \\
    \quad warning                           & 96 = 28.5\,\% \\
    \quad fail                              & 87 = 25.8\,\% \\
    Cases the Auditor could not audit       & 6 \\
    Structural violations                   & 0 \\
    \midrule
    Time: Axis A / B / C                    & 37\,s / 35.1\,h / 6.3\,h \\
    \bottomrule
  \end{tabular}
\end{table}

\begin{table}[t]
  \centering
  \caption{RAG pipeline: state of the documented runs.}
  \label{tab:agentic-baseline}
  \begin{tabular}{lr}
    \toprule
    Measure & Value \\
    \midrule
    Chunks produced (three successive runs) & 843 / 1,175 / 1,202 \\
    Normative chunks in the indexed run     & 466 \\
    Unit tests (all external calls mocked)  & 23 \\
    Test plans produced by these runs       & 0 \\
    \bottomrule
  \end{tabular}
\end{table}

\begin{table}[t]
  \centering
  \caption{RAG pipeline: measured diagnosis on the 1,202-chunk run.}
  \label{tab:diagnosis}
  \small\setlength{\tabcolsep}{4pt}
  \begin{tabular}{lrl}
    \toprule
    Measurement & Value & Consequence \\
    \midrule
    Distinct section labels                  & 25          & about 200 expected \\
    Chunks labelled with one ASN.1 OID       & 645 (54\,\%) & headings matched on OIDs \\
    Chunks labelled \texttt{8.1.2.2}         & 327 (27\,\%) & no heading for 40 pages \\
    Normative chunks holding a full clause  & 1 / 466     & chunking inactive \\
    Chunks with defined terms                & 0 / 1,202   & term injection inactive \\
    Cross-references resolved                & 22, all to OIDs & resolution inactive \\
    Planner clusters                         & 5           & feature sets meaningless \\
    \bottomrule
  \end{tabular}
\end{table}

\begin{table}[t]
  \centering
  \caption{RAG pipeline: extraction back-ends on ISO/IEC DIS 18013-5.}
  \label{tab:backends}
  \small\setlength{\tabcolsep}{4pt}
  \begin{tabular}{llllr}
    \toprule
    Back-end & Headings recovered & Ligatures & Font data & Time \\
    \midrule
    pdfplumber & 0.89 recall & broken & no  & 60\,s \\
    MarkItDown & none                         & broken & no  & 47\,s \\
    PyMuPDF    & 47 / 47 bookmarks     & intact & yes & 5\,s \\
    \bottomrule
  \end{tabular}
\end{table}

\begin{table}[t]
  \centering
  \caption{RAG pipeline: planned evaluation. Coverage is per normative
    chunk, the section labels being unreliable; the Axis~C rates are over the
    337 generated cases.}
  \label{tab:agentic-pending}
  \small\setlength{\tabcolsep}{4pt}
  \begin{tabular}{llrr}
    \toprule
    Axis & Metric & Target & Result \\
    \midrule
    A & Retrieval precision@5           &             & \tbd \\
    A & Retrieval recall@5              &             & \tbd \\
    A & Retrieval F1@5                  &             & \tbd \\
    \midrule
    B & Test cases generated            &             & 337 \\
    B & Normative chunks covered        & $\geq 80\,\%$ & 71.4\,\% \\
    B & Human score (5 criteria, 1--5)  & $\geq 3.5$    & \tbd \\
    \midrule
    C & Hallucination rate              & $< 10\,\%$    & 23.4\,\% \\
    C & Contradiction rate              & $< 5\,\%$     & 4.7\,\% \\
    C & Omission rate                   & $< 20\,\%$    & 43.3\,\% \\
    C & Adversarial detection rate      & $\geq 80\,\%$ & \tbd \\
    C & Auditor--human agreement (Cohen's $\kappa$) &   & \tbd \\
    \bottomrule
  \end{tabular}
\end{table}

% ---------------------------------------------------------------------------
\subsection{Method 2: deterministic pipeline}
\label{sec:exp-deterministic}

The second method is a fixed pipeline with no embeddings, no retrieval and no
agent framework: the control flow is written in code, the language model is
called only to extract and to write, and every reply is checked by rule.
Whatever fails a check is recorded as an explicit gap instead of being
emitted. Three successive versions of this pipeline were evaluated; they are
reported separately because each changes what the generator enumerates.

\subsubsection{Version 1: sentence-level generation}
\label{sec:exp-sentence}

\paragraph{Method.}
A rule-based ingest segments the document into clauses from its layout;
candidate normative sentences are selected by rule and one LLM call per clause
extracts requirements, each accepted only if its quote is verbatim and its
modality matches; one LLM call per requirement then writes a
nominal case and its negative and boundary variants, checked by a rule-based
verifier with a bounded retry. The model is \texttt{gemma4:31b}.

\paragraph{Results.}
Table~\ref{tab:sentence} gives the scores against the official suite.
Coverage is
4.2\,\%, and the three most-tested clauses score 3\,\%, 1\,\% and 0\,\%. The
generated cases are nevertheless mostly valid, which places the problem in
what is generated rather than in how it is written: the ingest recovered only
117 clauses, and one case per prose sentence does not have the granularity of
the official suite.

\begin{table}[t]
  \centering
  \caption{Deterministic pipeline, version 1 (\texttt{gemma4:31b}).}
  \label{tab:sentence}
  \begin{tabular}{lr}
    \toprule
    Measure & Value \\
    \midrule
    Clauses ingested                 & 117 \\
    Requirements extracted           & 215 \\
    Test cases generated             & 428 \\
    \quad nominal / negative / boundary & 215 / 186 / 27 \\
    \midrule
    Judge coverage, all citing cases & 29 / 687 = 4.2\,\% \\
    \quad Appendix 1                 & 3 / 14 \\
    \quad Appendix 2 (mDL)           & 16 / 451 \\
    \quad Appendix 3 (mDL reader)    & 10 / 222 \\
    \midrule
    Precision: grounded              & 52 / 60 \\
    Precision: actionable            & 50 / 60 \\
    Precision: expected result correct & 51 / 60 \\
    \bottomrule
  \end{tabular}
\end{table}

\subsubsection{Version 2: explicit test model and coverage items}
\label{sec:exp-testmodel}

\paragraph{Method.}
The official suite is an enumeration of objects against a fixed ladder of
checks, so an explicit test model was inserted between extraction and
generation. A second LLM stage extracts the objects the standard defines
(kind, presence, type, value domain, size, relations, each with a verbatim
quote); a deterministic stage derives coverage items from requirements and
objects by equivalence partitioning and boundary analysis; and one LLM call
per coverage item writes the test case, the kind of case being set by the tool
and never by the model. The ingest was extended to keep table rows as records,
and extraction accepts drafts individually instead of rejecting a whole batch.
The model is unchanged.

\paragraph{Results.}
Table~\ref{tab:testmodel} compares this run with version~1 on the same
model. Coverage of the official suite rises from 4.2\,\% to 13.8\,\%, the
best clause reaching 25.4\,\%, and 99.2\,\% of the coverage items produce a
case. Actionability does not move and becomes the quality ceiling. The
remaining limit is measured: twelve of the sixteen object-extraction failures
are truncated replies on the largest table clauses, which hold 361 table rows
and yielded 4 objects.

\begin{table}[t]
  \centering
  \caption{Deterministic pipeline, version 2 against version 1, same model
    (\texttt{gemma4:31b}).}
  \label{tab:testmodel}
  \begin{tabular}{lrr}
    \toprule
    Measure & Version 1 & Version 2 \\
    \midrule
    Clauses ingested            & 117 & 265 \\
    Requirements                & 215 & 842 \\
    Objects in the test model   & --  & 921 \\
    Coverage items              & --  & 3,257 \\
    Test cases                  & 428 & 3,231 \\
    Gaps                        & 14  & 66 \\
    \midrule
    Judge coverage              & 29 / 687 = 4.2\,\% & 95 / 687 = 13.8\,\% \\
    \quad Appendix 1            & 3 / 14   & 3 / 14 \\
    \quad Appendix 2 (mDL)      & 16 / 451 & 69 / 451 \\
    \quad Appendix 3 (mDL reader) & 10 / 222 & 23 / 222 \\
    \midrule
    Precision: grounded         & 52 / 60 & 56 / 60 \\
    Precision: actionable       & 50 / 60 & 50 / 60 \\
    Precision: expected result correct & 51 / 60 & 54 / 60 \\
    \bottomrule
  \end{tabular}
\end{table}

\subsubsection{Version 3: conformance-testing methodology}
\label{sec:exp-conformance}

\paragraph{Method.}
The third version adopts the methodology of the conformance-testing
standards the official suite follows (ISO/IEC 9646-2, ITU-T X.291). Four
elements are added to the pipeline. A deterministic grammar stage harvests
objects from formal declarations (CDDL, ABNF, JSON Schema, ASN.1). An
Implementation Conformance Statement (ICS) proforma is derived, with one item
per declared feature and a status of mandatory, optional or conditional, and
one \emph{capability} coverage item is emitted per mandatory or optional
item. Coverage items are organised in the four test groups of X.291
(capability, valid, invalid and inopportune behaviour), with the invalid side
of the ladder emitted for every object. Generation is batched per source. The
evaluation is rebuilt at the same time, with the judge-free structural matcher
and three explicit denominators. The full run of this version was made with
\texttt{qwen3.8:27b}, not with the model of version~2, so every
model-dependent figure below carries a change of model as well as a change of
method.

\paragraph{Results: the deterministic stages.}
The grammar, ICS and derivation stages use no model, so their output is
exactly reproducible from the input alone (Table~\ref{tab:deterministic}).
About a third of the 94 objects first harvested were malformed; the 152
harvested at the end are all well formed. The
share of due ICS items that carry a coverage item is 100\,\% by construction;
it is a soundness invariant of the derivation, not a score, and the
informative figure is the number of ICS items found.

\begin{table}[t]
  \centering
  \caption{Version 3: deterministic stages alone, no model, on ISO/IEC DIS
    18013-5 (265 clauses).}
  \label{tab:deterministic}
  \small\setlength{\tabcolsep}{4pt}
  \begin{tabular}{lrr}
    \toprule
    Measure & Start of development & End of development \\
    \midrule
    Objects harvested from formal notation & 94 & 152 \\
    ICS items                              & --   & 75 \\
    Coverage items                         & 250  & 604 \\
    \quad after deflation                  & --   & 596 \\
    Invalid-side coverage items            & 0    & 198 \\
    Test groups represented                & 2 of 4 & 4 of 4 \\
    \quad capability / valid / invalid     & --   & 75 / 331 / 190 \\
    Due ICS items with a coverage item     & --   & 75 / 75 \\
    \bottomrule
  \end{tabular}
\end{table}

\paragraph{Results: the full run.}
Table~\ref{tab:conformance-run} compares the full run with version~2.
Everything upstream of generation improved: the table clauses that had
truncated are now read, the object count rises by a factor of 1.73, and the
structural match on the comparable basis rises from 22.5\,\% to 34.9\,\%. The
generation stage, however, returned no test case in this run, so judge
coverage and precision are void, not low, and remain to be measured.

\begin{table}[t]
  \centering
  \caption{Version 3 against version 2. The two runs use different
    models. Groups: capability, valid, invalid and inopportune behaviour.
    G / A / C: grounded, actionable and correct.}
  \label{tab:conformance-run}
  \small\setlength{\tabcolsep}{4pt}
  \begin{tabular}{lrr}
    \toprule
    Measure & Version 2 & Version 3 \\
            & (\texttt{gemma4:31b}) & (\texttt{qwen3.8:27b}) \\
    \midrule
    Requirements                  & 842   & 810 \\
    Objects                       & 921   & 1,590 \\
    \quad of which from grammar   & --    & 152 \\
    ICS items                     & --    & 1,989 \\
    \quad mand. / opt. / cond. & -- & 1,438 / 394 / 157 \\
    Coverage items                & 3,257 & 7,978 \\
    \quad cap. / valid / invalid / inopp. & -- & 1,723 / 4,052 / 2,186 / 17 \\
    Due ICS items with a coverage item & -- & 1,832 / 1,832 \\
    \midrule
    Structural match, objects only & 168 / 747 = 22.5\,\% & 261 / 747 = 34.9\,\% \\
    Structural match, every source & 168 / 747 = 22.5\,\% & 276 / 747 = 36.9\,\% \\
    \midrule
    Test cases                    & 3,231 & 0 \\
    Judge coverage                & 95 / 687 = 13.8\,\% & void \\
    Precision (G / A / C, of 60) & 56 / 50 / 54 & void \\
    \bottomrule
  \end{tabular}
\end{table}

\paragraph{Results: denominators and macro averages.}
The size of the official suite is not its scored size. Of the 907 official
cases, 687 cite a clause of the standard under test, and only 250 cite a
clause identifier that exists in the 2020 draft used as input.
Table~\ref{tab:denominators} reports every rate against the three
denominators, with the macro averages over clauses and over test families
beside the per-case rate. On version~2, where both measures exist, the judge
and the structural matcher agree on 73.4\,\% of the cases both can score
(434 of 591).

\begin{table}[t]
  \centering
  \caption{Coverage of the official suite against three denominators. Micro is
    the per-case rate; macro is the mean of the per-clause or per-family
    rates. In version~2 no requirement carries a name, so the structural
    match is on named objects only.}
  \label{tab:denominators}
  \small\setlength{\tabcolsep}{4pt}
  \begin{tabular}{llrrrr}
    \toprule
    Scope & Run and measure & Scored & Micro & Macro & Macro \\
          &                 &        &       & clauses & families \\
    \midrule
    All official cases & V2, judge                      & 687 & 13.8\,\% & 17.0\,\% & 11.5\,\% \\
    ($n = 907$)        & V2, structural                 & 747 & 22.5\,\% & 20.5\,\% & 17.5\,\% \\
                       & V3, structural, objects only   & 747 & 34.9\,\% & 24.5\,\% & 39.0\,\% \\
                       & V3, structural, every source   & 747 & 36.9\,\% & 28.6\,\% & 40.6\,\% \\
    \midrule
    Citing the standard & V2, judge                     & 687 & 13.8\,\% & 17.0\,\% & 11.5\,\% \\
    ($n = 687$)        & V2, structural                 & 591 & 22.8\,\% & 20.5\,\% & 17.1\,\% \\
                       & V3, structural, objects only   & 591 & 34.7\,\% & 24.2\,\% & 39.9\,\% \\
                       & V3, structural, every source   & 591 & 36.7\,\% & 28.4\,\% & 40.9\,\% \\
    \midrule
    In-edition         & V2, judge                      & 250 & 16.8\,\% & 14.9\,\% & 23.3\,\% \\
    ($n = 250$)        & V2, structural                 & 207 & 18.8\,\% & 10.0\,\% & 14.8\,\% \\
                       & V3, structural, objects only   & 207 & 33.3\,\% & 20.9\,\% & 27.3\,\% \\
                       & V3, structural, every source   & 207 & 35.7\,\% & 32.5\,\% & 33.9\,\% \\
    \bottomrule
  \end{tabular}
\end{table}

\paragraph{Results: by check family and by clause.}
The gain is concentrated where the diagnosis of version~2 placed the limit
(Tables~\ref{tab:families} and~\ref{tab:clauses}): the encoding, consistency
and presence families, and the data-element clauses 7.2 and 7.2.1, which go
together from 40 to 118 matched cases of 166. Behaviour, the largest family of
the suite, stays at 16.4\,\%, size falls to zero, and clauses 9.2.2 and 9.3.2
are not matched at all in either run. In total 145 official cases became
matched and 37 stopped being matched.

\begin{table}[t]
  \centering
  \caption{Structural match by check family, matched of scorable official
    cases.}
  \label{tab:families}
  \begin{tabular}{lrrrr}
    \toprule
    Family & \multicolumn{2}{c}{Version 2} & \multicolumn{2}{c}{Version 3} \\
    \midrule
    Encoding     & 35 / 117 & 29.9\,\% & 85 / 117 & 72.6\,\% \\
    Consistency  & 7 / 31   & 22.6\,\% & 21 / 31  & 67.7\,\% \\
    Presence     & 45 / 137 & 32.8\,\% & 69 / 137 & 50.4\,\% \\
    Value domain & 56 / 165 & 33.9\,\% & 62 / 165 & 37.6\,\% \\
    Behaviour    & 20 / 220 & 9.1\,\%  & 36 / 220 & 16.4\,\% \\
    Reuse        & 2 / 24   & 8.3\,\%  & 3 / 24   & 12.5\,\% \\
    Size         & 3 / 53   & 5.7\,\%  & 0 / 53   & 0.0\,\% \\
    \midrule
    All scorable & 168 / 747 & 22.5\,\% & 276 / 747 & 36.9\,\% \\
    \bottomrule
  \end{tabular}
\end{table}

\begin{table}[t]
  \centering
  \caption{Structural match on the twelve most-tested clauses (2021
    numbering), matched of scorable official cases.}
  \label{tab:clauses}
  \begin{tabular}{lrr}
    \toprule
    Clause & Version 2 & Version 3 \\
    \midrule
    7.2.1       & 16 / 90 & 69 / 90 \\
    7.2         & 24 / 76 & 49 / 76 \\
    9.1.2.5     & 9 / 53  & 17 / 53 \\
    9.1.2.4     & 27 / 82 & 30 / 82 \\
    8.3.2.1.2.2 & 19 / 86 & 22 / 86 \\
    8.2.1.1     & 12 / 49 & 15 / 49 \\
    9.1.4       & 6 / 34  & 9 / 34 \\
    8.3.2.1.2   & 5 / 27  & 6 / 27 \\
    8.3.2.1.2.3 & 9 / 32  & 9 / 32 \\
    9.1.1.4     & 13 / 54 & 8 / 54 \\
    9.2.2       & 0 / 33  & 0 / 33 \\
    9.3.2       & 0 / 30  & 0 / 30 \\
    \bottomrule
  \end{tabular}
\end{table}

% ---------------------------------------------------------------------------
\subsection{Method 3: agentic harness}
\label{sec:exp-harness}

% TODO: complete this subsection once the runs are done.
\paragraph{Method.}
% TODO: describe the protocol. Points to fix: the prompt given to the harness,
% what it may read (the standard only, never TS 18013-6), the output format
% (the same plan schema, so that the same evaluation scripts apply), the
% models, and the judge model, which should be the same for every row.
The third method gives the task to an industrial agentic coding harness,
Claude Code, with no purpose-built pipeline around it. The harness receives
the same input document and is asked to produce a test plan in the same
format, which is then scored with the same evaluation as the deterministic
pipeline. Two families of models are run through the harness: Anthropic
models (\tbd) and open-source models (\tbd).

\paragraph{Results.}
Table~\ref{tab:harness} gives the results per model. \tbd

\begin{table}[t]
  \centering
  \caption{Agentic harness: results per model (pending). Judge: judge
    coverage; Struct.: structural match; G / A / C: grounded, actionable and
    correct, of 60 sampled cases.}
  \label{tab:harness}
  \small\setlength{\tabcolsep}{4pt}
  \begin{tabular}{llrrrrr}
    \toprule
    Harness & Model & Cases & Judge & Struct. & G / A / C & Time \\
    \midrule
    Claude Code & Anthropic, \tbd   & \tbd & \tbd & \tbd & \tbd & \tbd \\
    Claude Code & Anthropic, \tbd   & \tbd & \tbd & \tbd & \tbd & \tbd \\
    Claude Code & Open source, \tbd & \tbd & \tbd & \tbd & \tbd & \tbd \\
    Claude Code & Open source, \tbd & \tbd & \tbd & \tbd & \tbd & \tbd \\
    \bottomrule
  \end{tabular}
\end{table}

% ---------------------------------------------------------------------------
\subsection{Comparison of the three methods}
\label{sec:exp-comparison}

% TODO: write the comparison once the RAG and harness results are in.
% Note: the structural match needs a plan that carries coverage items with a
% named source object; decide how it applies to the RAG and harness outputs.
Table~\ref{tab:comparison} places the three methods side by side on the
metrics of Section~\ref{sec:exp-setting}. \tbd

\begin{table}[t]
  \centering
  \caption{The three methods on ISO/IEC DIS 18013-5 (harness results and
    the scoring of the RAG plan pending). Judge: judge coverage, of the 687 official cases citing the
    standard; Struct.: structural match, of the 747 scorable official cases;
    G / A / C: grounded, actionable and correct, of 60 sampled cases.}
  \label{tab:comparison}
  \small\setlength{\tabcolsep}{4pt}
  \begin{tabular}{llrrrrr}
    \toprule
    Method & Model & Cases & Judge & Struct. & G / A / C & Time \\
    \midrule
    1. RAG pipeline       & \texttt{gemma4:31b}  & 337   & \tbd     & \tbd     & \tbd         & 41.4\,h \\
    \midrule
    2. Deterministic, v1  & \texttt{gemma4:31b}  & 428   & 4.2\,\%  & --       & 52 / 50 / 51 & 0.75\,h \\
    2. Deterministic, v2  & \texttt{gemma4:31b}  & 3,231 & 13.8\,\% & 22.5\,\% & 56 / 50 / 54 & 12.7\,h \\
    2. Deterministic, v3  & \texttt{qwen3.8:27b} & 0     & void     & 36.9\,\% & void         & 126.1\,h \\
    \midrule
    3. Harness            & Anthropic, \tbd      & \tbd  & \tbd     & \tbd     & \tbd         & \tbd \\
    3. Harness            & Open source, \tbd    & \tbd  & \tbd     & \tbd     & \tbd         & \tbd \\
    \bottomrule
  \end{tabular}
\end{table}
 from the main article. Requires \usepackage{booktabs}.
% Every \tbd marks a result that has not been computed yet.
% Sources of the figures: docs/article/outline.md and, under method2-deterministic/,
% docs/evaluation/*.md, eval/eval_remote_1/eval/summary.txt, eval/eval_remote_2/eval/summary.txt.
% Method 1 full run: method1-rag/data/output/vps_run_2026-11-19/ (plan, audit report, run_all.log).
\providecommand{\tbd}{\textit{tbd}}

\section{Experiments and Results}
\label{sec:experiments}

Three methods are compared on the same task, generating a conformance test
plan from the text of a standard: a retrieval-augmented agentic pipeline
(Section~\ref{sec:exp-rag}), a deterministic pipeline in which the language
model only fills the steps a fixed procedure assigns to it
(Section~\ref{sec:exp-deterministic}), and an industrial agentic harness
given the task directly (Section~\ref{sec:exp-harness}). Each subsection
gives the method in a short paragraph and then its results;
Section~\ref{sec:exp-comparison} places the three side by side.

\subsection{Experimental setting}
\label{sec:exp-setting}

\paragraph{Input and ground truth.}
All experiments use the same input, the 126-page draft ISO/IEC DIS
18013-5:2020, and the same ground truth, the three appendices of ISO/IEC TS
18013-6:2025, parsed by rule into structured test cases
(Table~\ref{tab:ground-truth}). The 687 official cases that cite the standard
refer to 54 distinct clauses of the published 2021 edition, whereas the input
is the 2020 draft; this edition mismatch is a confound in every coverage
figure below and is quantified in Section~\ref{sec:exp-conformance}.

\begin{table}[t]
  \centering
  \caption{Ground truth: official test cases of ISO/IEC TS 18013-6:2025.}
  \label{tab:ground-truth}
  \begin{tabular}{lrr}
    \toprule
    Appendix & Test cases & Citing an 18013-5 clause \\
    \midrule
    1 -- Common and certificates & 138 & 14 \\
    2 -- mDL                     & 460 & 451 \\
    3 -- mDL reader              & 309 & 222 \\
    \midrule
    Total                        & 907 & 687 \\
    \bottomrule
  \end{tabular}
\end{table}

\paragraph{Metrics.}
\emph{Judge coverage} is the share of official cases for which an LLM judge,
shown the five generated cases closest to the official one, finds at least one
that verifies the same check (same data item or behaviour, same condition,
same expected outcome). \emph{Structural match} uses no model: each official
case is reduced by rule to a pair (data item, check type), and it is matched
when the generated plan carries a coverage item whose source object is named
by that data item and whose check belongs to the same family; 747 of the 907
official cases can be reduced to such a pair and are called \emph{scorable}.
\emph{Precision} is judged on a fixed random sample of 60 generated cases
against three criteria: the case is \emph{grounded} in its requirement,
\emph{actionable} by a test engineer, and its expected result is
\emph{correct}. The judge belongs to the same model family as the generator,
which is a known self-preference risk; the structural matcher was introduced
to remove it from the headline figure.

% ---------------------------------------------------------------------------
\subsection{Method 1: RAG pipeline}
\label{sec:exp-rag}

\paragraph{Method.}
The first method is an agentic retrieval-augmented pipeline organised in
three stages. Axis~A extracts the PDF, splits it with a structure-aware
two-pass chunker and indexes the chunks for hybrid dense and sparse retrieval
fused by reciprocal rank. Axis~B runs a Planner and then a chain of three
agents (Analyst, RAG Router, Designer) once per normative chunk, with
schema-constrained generation, an error-reflective retry and a BERTScore
self-evaluation gate, and assembles the output as an ISO/IEC/IEEE 29119-3
test plan. Axis~C audits each generated case with structural guardrails, an
embedding-distance pre-audit and an Auditor agent whose verdict is overridden
by a deterministic rule. The pipeline was developed with \texttt{llama3.2};
the full run reported here uses \texttt{gemma4:31b} for every agent, served
locally with a context of 4096 tokens.

\paragraph{Results: intermediate artefacts.}
The first documented runs produced chunks and an index but no test plan
(Table~\ref{tab:agentic-baseline}). The measurements of
Table~\ref{tab:diagnosis} show that a single defect, a heading detector that
matched object identifiers in the annexes instead of clause numbers, silently
disabled the chunking, cross-reference, defined-term and planning mechanisms.
A comparison of extraction back-ends (Table~\ref{tab:backends}) located the
cause in the extraction layer.

\paragraph{Results: the full run.}
One full run of the three axes was then made with the heading defect still
present (Table~\ref{tab:agentic-run}). Axis~B wrote a test case for 337 of
the 472 normative chunks (71.4\,\%); the other 135 were dropped because an
agent did not return valid JSON within its retries, the RAG Router alone
accounting for 102 of them. The loss does not depend on the strength of the
requirement (72\,\% of the chunks with a \emph{shall}, 69\,\% with a
\emph{should}, 71\,\% with a \emph{may}) but on the place in the document: it
is highest on the certificate-profile tables. Axis~C passed 154 of the 337
cases (45.7\,\%). Its hallucination and omission rates miss their targets and
its contradiction rate meets it (Table~\ref{tab:agentic-pending}). These rates
are an optimistic estimate: the Auditor is the model that wrote the cases, and
the six cases it failed to audit are counted as free of findings. The schema
check found no violation, and the embedding-distance pre-audit flagged no
case, its score being the same for passed and failed cases (0.84).

The heading defect is visible in the plan. The 1,246 chunks carry 27 distinct
section labels, so the 337 cases share 97 identifiers, 127 of them a single
one, the five feature sets have no meaningful name, and a case can only be
traced to its chunk, not to a clause. For that reason coverage is reported
per normative chunk and not per section, and the run has not been scored
against the official suite: both the judge and the structural matcher need a
clause or a named object on the generated side. The cases are also short,
261 of 337 having exactly two steps.

\begin{table}[t]
  \centering
  \caption{RAG pipeline: full run on ISO/IEC DIS 18013-5
    (\texttt{gemma4:31b}).}
  \label{tab:agentic-run}
  \begin{tabular}{lr}
    \toprule
    Measure & Value \\
    \midrule
    Chunks indexed (Axis A)                 & 1,246 \\
    \quad normative                         & 472 \\
    Distinct section labels                 & 27 \\
    \midrule
    Test cases generated (Axis B)           & 337 / 472 = 71.4\,\% \\
    \quad \emph{shall} / \emph{should} / \emph{may} & 248 / 61 / 28 \\
    Chunks dropped                          & 135 \\
    \quad RAG Router / Designer / Analyst   & 102 / 31 / 2 \\
    Distinct test case identifiers          & 97 \\
    Steps per case, mean                    & 2.3 \\
    \midrule
    Audit verdict (Axis C): pass            & 154 = 45.7\,\% \\
    \quad warning                           & 96 = 28.5\,\% \\
    \quad fail                              & 87 = 25.8\,\% \\
    Cases the Auditor could not audit       & 6 \\
    Structural violations                   & 0 \\
    \midrule
    Time: Axis A / B / C                    & 37\,s / 35.1\,h / 6.3\,h \\
    \bottomrule
  \end{tabular}
\end{table}

\begin{table}[t]
  \centering
  \caption{RAG pipeline: state of the documented runs.}
  \label{tab:agentic-baseline}
  \begin{tabular}{lr}
    \toprule
    Measure & Value \\
    \midrule
    Chunks produced (three successive runs) & 843 / 1,175 / 1,202 \\
    Normative chunks in the indexed run     & 466 \\
    Unit tests (all external calls mocked)  & 23 \\
    Test plans produced by these runs       & 0 \\
    \bottomrule
  \end{tabular}
\end{table}

\begin{table}[t]
  \centering
  \caption{RAG pipeline: measured diagnosis on the 1,202-chunk run.}
  \label{tab:diagnosis}
  \small\setlength{\tabcolsep}{4pt}
  \begin{tabular}{lrl}
    \toprule
    Measurement & Value & Consequence \\
    \midrule
    Distinct section labels                  & 25          & about 200 expected \\
    Chunks labelled with one ASN.1 OID       & 645 (54\,\%) & headings matched on OIDs \\
    Chunks labelled \texttt{8.1.2.2}         & 327 (27\,\%) & no heading for 40 pages \\
    Normative chunks holding a full clause  & 1 / 466     & chunking inactive \\
    Chunks with defined terms                & 0 / 1,202   & term injection inactive \\
    Cross-references resolved                & 22, all to OIDs & resolution inactive \\
    Planner clusters                         & 5           & feature sets meaningless \\
    \bottomrule
  \end{tabular}
\end{table}

\begin{table}[t]
  \centering
  \caption{RAG pipeline: extraction back-ends on ISO/IEC DIS 18013-5.}
  \label{tab:backends}
  \small\setlength{\tabcolsep}{4pt}
  \begin{tabular}{llllr}
    \toprule
    Back-end & Headings recovered & Ligatures & Font data & Time \\
    \midrule
    pdfplumber & 0.89 recall & broken & no  & 60\,s \\
    MarkItDown & none                         & broken & no  & 47\,s \\
    PyMuPDF    & 47 / 47 bookmarks     & intact & yes & 5\,s \\
    \bottomrule
  \end{tabular}
\end{table}

\begin{table}[t]
  \centering
  \caption{RAG pipeline: planned evaluation. Coverage is per normative
    chunk, the section labels being unreliable; the Axis~C rates are over the
    337 generated cases.}
  \label{tab:agentic-pending}
  \small\setlength{\tabcolsep}{4pt}
  \begin{tabular}{llrr}
    \toprule
    Axis & Metric & Target & Result \\
    \midrule
    A & Retrieval precision@5           &             & \tbd \\
    A & Retrieval recall@5              &             & \tbd \\
    A & Retrieval F1@5                  &             & \tbd \\
    \midrule
    B & Test cases generated            &             & 337 \\
    B & Normative chunks covered        & $\geq 80\,\%$ & 71.4\,\% \\
    B & Human score (5 criteria, 1--5)  & $\geq 3.5$    & \tbd \\
    \midrule
    C & Hallucination rate              & $< 10\,\%$    & 23.4\,\% \\
    C & Contradiction rate              & $< 5\,\%$     & 4.7\,\% \\
    C & Omission rate                   & $< 20\,\%$    & 43.3\,\% \\
    C & Adversarial detection rate      & $\geq 80\,\%$ & \tbd \\
    C & Auditor--human agreement (Cohen's $\kappa$) &   & \tbd \\
    \bottomrule
  \end{tabular}
\end{table}

% ---------------------------------------------------------------------------
\subsection{Method 2: deterministic pipeline}
\label{sec:exp-deterministic}

The second method is a fixed pipeline with no embeddings, no retrieval and no
agent framework: the control flow is written in code, the language model is
called only to extract and to write, and every reply is checked by rule.
Whatever fails a check is recorded as an explicit gap instead of being
emitted. Three successive versions of this pipeline were evaluated; they are
reported separately because each changes what the generator enumerates.

\subsubsection{Version 1: sentence-level generation}
\label{sec:exp-sentence}

\paragraph{Method.}
A rule-based ingest segments the document into clauses from its layout;
candidate normative sentences are selected by rule and one LLM call per clause
extracts requirements, each accepted only if its quote is verbatim and its
modality matches; one LLM call per requirement then writes a
nominal case and its negative and boundary variants, checked by a rule-based
verifier with a bounded retry. The model is \texttt{gemma4:31b}.

\paragraph{Results.}
Table~\ref{tab:sentence} gives the scores against the official suite.
Coverage is
4.2\,\%, and the three most-tested clauses score 3\,\%, 1\,\% and 0\,\%. The
generated cases are nevertheless mostly valid, which places the problem in
what is generated rather than in how it is written: the ingest recovered only
117 clauses, and one case per prose sentence does not have the granularity of
the official suite.

\begin{table}[t]
  \centering
  \caption{Deterministic pipeline, version 1 (\texttt{gemma4:31b}).}
  \label{tab:sentence}
  \begin{tabular}{lr}
    \toprule
    Measure & Value \\
    \midrule
    Clauses ingested                 & 117 \\
    Requirements extracted           & 215 \\
    Test cases generated             & 428 \\
    \quad nominal / negative / boundary & 215 / 186 / 27 \\
    \midrule
    Judge coverage, all citing cases & 29 / 687 = 4.2\,\% \\
    \quad Appendix 1                 & 3 / 14 \\
    \quad Appendix 2 (mDL)           & 16 / 451 \\
    \quad Appendix 3 (mDL reader)    & 10 / 222 \\
    \midrule
    Precision: grounded              & 52 / 60 \\
    Precision: actionable            & 50 / 60 \\
    Precision: expected result correct & 51 / 60 \\
    \bottomrule
  \end{tabular}
\end{table}

\subsubsection{Version 2: explicit test model and coverage items}
\label{sec:exp-testmodel}

\paragraph{Method.}
The official suite is an enumeration of objects against a fixed ladder of
checks, so an explicit test model was inserted between extraction and
generation. A second LLM stage extracts the objects the standard defines
(kind, presence, type, value domain, size, relations, each with a verbatim
quote); a deterministic stage derives coverage items from requirements and
objects by equivalence partitioning and boundary analysis; and one LLM call
per coverage item writes the test case, the kind of case being set by the tool
and never by the model. The ingest was extended to keep table rows as records,
and extraction accepts drafts individually instead of rejecting a whole batch.
The model is unchanged.

\paragraph{Results.}
Table~\ref{tab:testmodel} compares this run with version~1 on the same
model. Coverage of the official suite rises from 4.2\,\% to 13.8\,\%, the
best clause reaching 25.4\,\%, and 99.2\,\% of the coverage items produce a
case. Actionability does not move and becomes the quality ceiling. The
remaining limit is measured: twelve of the sixteen object-extraction failures
are truncated replies on the largest table clauses, which hold 361 table rows
and yielded 4 objects.

\begin{table}[t]
  \centering
  \caption{Deterministic pipeline, version 2 against version 1, same model
    (\texttt{gemma4:31b}).}
  \label{tab:testmodel}
  \begin{tabular}{lrr}
    \toprule
    Measure & Version 1 & Version 2 \\
    \midrule
    Clauses ingested            & 117 & 265 \\
    Requirements                & 215 & 842 \\
    Objects in the test model   & --  & 921 \\
    Coverage items              & --  & 3,257 \\
    Test cases                  & 428 & 3,231 \\
    Gaps                        & 14  & 66 \\
    \midrule
    Judge coverage              & 29 / 687 = 4.2\,\% & 95 / 687 = 13.8\,\% \\
    \quad Appendix 1            & 3 / 14   & 3 / 14 \\
    \quad Appendix 2 (mDL)      & 16 / 451 & 69 / 451 \\
    \quad Appendix 3 (mDL reader) & 10 / 222 & 23 / 222 \\
    \midrule
    Precision: grounded         & 52 / 60 & 56 / 60 \\
    Precision: actionable       & 50 / 60 & 50 / 60 \\
    Precision: expected result correct & 51 / 60 & 54 / 60 \\
    \bottomrule
  \end{tabular}
\end{table}

\subsubsection{Version 3: conformance-testing methodology}
\label{sec:exp-conformance}

\paragraph{Method.}
The third version adopts the methodology of the conformance-testing
standards the official suite follows (ISO/IEC 9646-2, ITU-T X.291). Four
elements are added to the pipeline. A deterministic grammar stage harvests
objects from formal declarations (CDDL, ABNF, JSON Schema, ASN.1). An
Implementation Conformance Statement (ICS) proforma is derived, with one item
per declared feature and a status of mandatory, optional or conditional, and
one \emph{capability} coverage item is emitted per mandatory or optional
item. Coverage items are organised in the four test groups of X.291
(capability, valid, invalid and inopportune behaviour), with the invalid side
of the ladder emitted for every object. Generation is batched per source. The
evaluation is rebuilt at the same time, with the judge-free structural matcher
and three explicit denominators. The full run of this version was made with
\texttt{qwen3.8:27b}, not with the model of version~2, so every
model-dependent figure below carries a change of model as well as a change of
method.

\paragraph{Results: the deterministic stages.}
The grammar, ICS and derivation stages use no model, so their output is
exactly reproducible from the input alone (Table~\ref{tab:deterministic}).
About a third of the 94 objects first harvested were malformed; the 152
harvested at the end are all well formed. The
share of due ICS items that carry a coverage item is 100\,\% by construction;
it is a soundness invariant of the derivation, not a score, and the
informative figure is the number of ICS items found.

\begin{table}[t]
  \centering
  \caption{Version 3: deterministic stages alone, no model, on ISO/IEC DIS
    18013-5 (265 clauses).}
  \label{tab:deterministic}
  \small\setlength{\tabcolsep}{4pt}
  \begin{tabular}{lrr}
    \toprule
    Measure & Start of development & End of development \\
    \midrule
    Objects harvested from formal notation & 94 & 152 \\
    ICS items                              & --   & 75 \\
    Coverage items                         & 250  & 604 \\
    \quad after deflation                  & --   & 596 \\
    Invalid-side coverage items            & 0    & 198 \\
    Test groups represented                & 2 of 4 & 4 of 4 \\
    \quad capability / valid / invalid     & --   & 75 / 331 / 190 \\
    Due ICS items with a coverage item     & --   & 75 / 75 \\
    \bottomrule
  \end{tabular}
\end{table}

\paragraph{Results: the full run.}
Table~\ref{tab:conformance-run} compares the full run with version~2.
Everything upstream of generation improved: the table clauses that had
truncated are now read, the object count rises by a factor of 1.73, and the
structural match on the comparable basis rises from 22.5\,\% to 34.9\,\%. The
generation stage, however, returned no test case in this run, so judge
coverage and precision are void, not low, and remain to be measured.

\begin{table}[t]
  \centering
  \caption{Version 3 against version 2. The two runs use different
    models. Groups: capability, valid, invalid and inopportune behaviour.
    G / A / C: grounded, actionable and correct.}
  \label{tab:conformance-run}
  \small\setlength{\tabcolsep}{4pt}
  \begin{tabular}{lrr}
    \toprule
    Measure & Version 2 & Version 3 \\
            & (\texttt{gemma4:31b}) & (\texttt{qwen3.8:27b}) \\
    \midrule
    Requirements                  & 842   & 810 \\
    Objects                       & 921   & 1,590 \\
    \quad of which from grammar   & --    & 152 \\
    ICS items                     & --    & 1,989 \\
    \quad mand. / opt. / cond. & -- & 1,438 / 394 / 157 \\
    Coverage items                & 3,257 & 7,978 \\
    \quad cap. / valid / invalid / inopp. & -- & 1,723 / 4,052 / 2,186 / 17 \\
    Due ICS items with a coverage item & -- & 1,832 / 1,832 \\
    \midrule
    Structural match, objects only & 168 / 747 = 22.5\,\% & 261 / 747 = 34.9\,\% \\
    Structural match, every source & 168 / 747 = 22.5\,\% & 276 / 747 = 36.9\,\% \\
    \midrule
    Test cases                    & 3,231 & 0 \\
    Judge coverage                & 95 / 687 = 13.8\,\% & void \\
    Precision (G / A / C, of 60) & 56 / 50 / 54 & void \\
    \bottomrule
  \end{tabular}
\end{table}

\paragraph{Results: denominators and macro averages.}
The size of the official suite is not its scored size. Of the 907 official
cases, 687 cite a clause of the standard under test, and only 250 cite a
clause identifier that exists in the 2020 draft used as input.
Table~\ref{tab:denominators} reports every rate against the three
denominators, with the macro averages over clauses and over test families
beside the per-case rate. On version~2, where both measures exist, the judge
and the structural matcher agree on 73.4\,\% of the cases both can score
(434 of 591).

\begin{table}[t]
  \centering
  \caption{Coverage of the official suite against three denominators. Micro is
    the per-case rate; macro is the mean of the per-clause or per-family
    rates. In version~2 no requirement carries a name, so the structural
    match is on named objects only.}
  \label{tab:denominators}
  \small\setlength{\tabcolsep}{4pt}
  \begin{tabular}{llrrrr}
    \toprule
    Scope & Run and measure & Scored & Micro & Macro & Macro \\
          &                 &        &       & clauses & families \\
    \midrule
    All official cases & V2, judge                      & 687 & 13.8\,\% & 17.0\,\% & 11.5\,\% \\
    ($n = 907$)        & V2, structural                 & 747 & 22.5\,\% & 20.5\,\% & 17.5\,\% \\
                       & V3, structural, objects only   & 747 & 34.9\,\% & 24.5\,\% & 39.0\,\% \\
                       & V3, structural, every source   & 747 & 36.9\,\% & 28.6\,\% & 40.6\,\% \\
    \midrule
    Citing the standard & V2, judge                     & 687 & 13.8\,\% & 17.0\,\% & 11.5\,\% \\
    ($n = 687$)        & V2, structural                 & 591 & 22.8\,\% & 20.5\,\% & 17.1\,\% \\
                       & V3, structural, objects only   & 591 & 34.7\,\% & 24.2\,\% & 39.9\,\% \\
                       & V3, structural, every source   & 591 & 36.7\,\% & 28.4\,\% & 40.9\,\% \\
    \midrule
    In-edition         & V2, judge                      & 250 & 16.8\,\% & 14.9\,\% & 23.3\,\% \\
    ($n = 250$)        & V2, structural                 & 207 & 18.8\,\% & 10.0\,\% & 14.8\,\% \\
                       & V3, structural, objects only   & 207 & 33.3\,\% & 20.9\,\% & 27.3\,\% \\
                       & V3, structural, every source   & 207 & 35.7\,\% & 32.5\,\% & 33.9\,\% \\
    \bottomrule
  \end{tabular}
\end{table}

\paragraph{Results: by check family and by clause.}
The gain is concentrated where the diagnosis of version~2 placed the limit
(Tables~\ref{tab:families} and~\ref{tab:clauses}): the encoding, consistency
and presence families, and the data-element clauses 7.2 and 7.2.1, which go
together from 40 to 118 matched cases of 166. Behaviour, the largest family of
the suite, stays at 16.4\,\%, size falls to zero, and clauses 9.2.2 and 9.3.2
are not matched at all in either run. In total 145 official cases became
matched and 37 stopped being matched.

\begin{table}[t]
  \centering
  \caption{Structural match by check family, matched of scorable official
    cases.}
  \label{tab:families}
  \begin{tabular}{lrrrr}
    \toprule
    Family & \multicolumn{2}{c}{Version 2} & \multicolumn{2}{c}{Version 3} \\
    \midrule
    Encoding     & 35 / 117 & 29.9\,\% & 85 / 117 & 72.6\,\% \\
    Consistency  & 7 / 31   & 22.6\,\% & 21 / 31  & 67.7\,\% \\
    Presence     & 45 / 137 & 32.8\,\% & 69 / 137 & 50.4\,\% \\
    Value domain & 56 / 165 & 33.9\,\% & 62 / 165 & 37.6\,\% \\
    Behaviour    & 20 / 220 & 9.1\,\%  & 36 / 220 & 16.4\,\% \\
    Reuse        & 2 / 24   & 8.3\,\%  & 3 / 24   & 12.5\,\% \\
    Size         & 3 / 53   & 5.7\,\%  & 0 / 53   & 0.0\,\% \\
    \midrule
    All scorable & 168 / 747 & 22.5\,\% & 276 / 747 & 36.9\,\% \\
    \bottomrule
  \end{tabular}
\end{table}

\begin{table}[t]
  \centering
  \caption{Structural match on the twelve most-tested clauses (2021
    numbering), matched of scorable official cases.}
  \label{tab:clauses}
  \begin{tabular}{lrr}
    \toprule
    Clause & Version 2 & Version 3 \\
    \midrule
    7.2.1       & 16 / 90 & 69 / 90 \\
    7.2         & 24 / 76 & 49 / 76 \\
    9.1.2.5     & 9 / 53  & 17 / 53 \\
    9.1.2.4     & 27 / 82 & 30 / 82 \\
    8.3.2.1.2.2 & 19 / 86 & 22 / 86 \\
    8.2.1.1     & 12 / 49 & 15 / 49 \\
    9.1.4       & 6 / 34  & 9 / 34 \\
    8.3.2.1.2   & 5 / 27  & 6 / 27 \\
    8.3.2.1.2.3 & 9 / 32  & 9 / 32 \\
    9.1.1.4     & 13 / 54 & 8 / 54 \\
    9.2.2       & 0 / 33  & 0 / 33 \\
    9.3.2       & 0 / 30  & 0 / 30 \\
    \bottomrule
  \end{tabular}
\end{table}

% ---------------------------------------------------------------------------
\subsection{Method 3: agentic harness}
\label{sec:exp-harness}

% TODO: complete this subsection once the runs are done.
\paragraph{Method.}
% TODO: describe the protocol. Points to fix: the prompt given to the harness,
% what it may read (the standard only, never TS 18013-6), the output format
% (the same plan schema, so that the same evaluation scripts apply), the
% models, and the judge model, which should be the same for every row.
The third method gives the task to an industrial agentic coding harness,
Claude Code, with no purpose-built pipeline around it. The harness receives
the same input document and is asked to produce a test plan in the same
format, which is then scored with the same evaluation as the deterministic
pipeline. Two families of models are run through the harness: Anthropic
models (\tbd) and open-source models (\tbd).

\paragraph{Results.}
Table~\ref{tab:harness} gives the results per model. \tbd

\begin{table}[t]
  \centering
  \caption{Agentic harness: results per model (pending). Judge: judge
    coverage; Struct.: structural match; G / A / C: grounded, actionable and
    correct, of 60 sampled cases.}
  \label{tab:harness}
  \small\setlength{\tabcolsep}{4pt}
  \begin{tabular}{llrrrrr}
    \toprule
    Harness & Model & Cases & Judge & Struct. & G / A / C & Time \\
    \midrule
    Claude Code & Anthropic, \tbd   & \tbd & \tbd & \tbd & \tbd & \tbd \\
    Claude Code & Anthropic, \tbd   & \tbd & \tbd & \tbd & \tbd & \tbd \\
    Claude Code & Open source, \tbd & \tbd & \tbd & \tbd & \tbd & \tbd \\
    Claude Code & Open source, \tbd & \tbd & \tbd & \tbd & \tbd & \tbd \\
    \bottomrule
  \end{tabular}
\end{table}

% ---------------------------------------------------------------------------
\subsection{Comparison of the three methods}
\label{sec:exp-comparison}

% TODO: write the comparison once the RAG and harness results are in.
% Note: the structural match needs a plan that carries coverage items with a
% named source object; decide how it applies to the RAG and harness outputs.
Table~\ref{tab:comparison} places the three methods side by side on the
metrics of Section~\ref{sec:exp-setting}. \tbd

\begin{table}[t]
  \centering
  \caption{The three methods on ISO/IEC DIS 18013-5 (harness results and
    the scoring of the RAG plan pending). Judge: judge coverage, of the 687 official cases citing the
    standard; Struct.: structural match, of the 747 scorable official cases;
    G / A / C: grounded, actionable and correct, of 60 sampled cases.}
  \label{tab:comparison}
  \small\setlength{\tabcolsep}{4pt}
  \begin{tabular}{llrrrrr}
    \toprule
    Method & Model & Cases & Judge & Struct. & G / A / C & Time \\
    \midrule
    1. RAG pipeline       & \texttt{gemma4:31b}  & 337   & \tbd     & \tbd     & \tbd         & 41.4\,h \\
    \midrule
    2. Deterministic, v1  & \texttt{gemma4:31b}  & 428   & 4.2\,\%  & --       & 52 / 50 / 51 & 0.75\,h \\
    2. Deterministic, v2  & \texttt{gemma4:31b}  & 3,231 & 13.8\,\% & 22.5\,\% & 56 / 50 / 54 & 12.7\,h \\
    2. Deterministic, v3  & \texttt{qwen3.8:27b} & 0     & void     & 36.9\,\% & void         & 126.1\,h \\
    \midrule
    3. Harness            & Anthropic, \tbd      & \tbd  & \tbd     & \tbd     & \tbd         & \tbd \\
    3. Harness            & Open source, \tbd    & \tbd  & \tbd     & \tbd     & \tbd         & \tbd \\
    \bottomrule
  \end{tabular}
\end{table}
 from the main article. Requires \usepackage{booktabs}.
% Every \tbd marks a result that has not been computed yet.
% Sources of the figures: docs/article/outline.md and, under method2-deterministic/,
% docs/evaluation/*.md, eval/eval_remote_1/eval/summary.txt, eval/eval_remote_2/eval/summary.txt.
% Method 1 full run: method1-rag/data/output/vps_run_2026-11-19/ (plan, audit report, run_all.log).
\providecommand{\tbd}{\textit{tbd}}

\section{Experiments and Results}
\label{sec:experiments}

Three methods are compared on the same task, generating a conformance test
plan from the text of a standard: a retrieval-augmented agentic pipeline
(Section~\ref{sec:exp-rag}), a deterministic pipeline in which the language
model only fills the steps a fixed procedure assigns to it
(Section~\ref{sec:exp-deterministic}), and an industrial agentic harness
given the task directly (Section~\ref{sec:exp-harness}). Each subsection
gives the method in a short paragraph and then its results;
Section~\ref{sec:exp-comparison} places the three side by side.

\subsection{Experimental setting}
\label{sec:exp-setting}

\paragraph{Input and ground truth.}
All experiments use the same input, the 126-page draft ISO/IEC DIS
18013-5:2020, and the same ground truth, the three appendices of ISO/IEC TS
18013-6:2025, parsed by rule into structured test cases
(Table~\ref{tab:ground-truth}). The 687 official cases that cite the standard
refer to 54 distinct clauses of the published 2021 edition, whereas the input
is the 2020 draft; this edition mismatch is a confound in every coverage
figure below and is quantified in Section~\ref{sec:exp-conformance}.

\begin{table}[t]
  \centering
  \caption{Ground truth: official test cases of ISO/IEC TS 18013-6:2025.}
  \label{tab:ground-truth}
  \begin{tabular}{lrr}
    \toprule
    Appendix & Test cases & Citing an 18013-5 clause \\
    \midrule
    1 -- Common and certificates & 138 & 14 \\
    2 -- mDL                     & 460 & 451 \\
    3 -- mDL reader              & 309 & 222 \\
    \midrule
    Total                        & 907 & 687 \\
    \bottomrule
  \end{tabular}
\end{table}

\paragraph{Metrics.}
\emph{Judge coverage} is the share of official cases for which an LLM judge,
shown the five generated cases closest to the official one, finds at least one
that verifies the same check (same data item or behaviour, same condition,
same expected outcome). \emph{Structural match} uses no model: each official
case is reduced by rule to a pair (data item, check type), and it is matched
when the generated plan carries a coverage item whose source object is named
by that data item and whose check belongs to the same family; 747 of the 907
official cases can be reduced to such a pair and are called \emph{scorable}.
\emph{Precision} is judged on a fixed random sample of 60 generated cases
against three criteria: the case is \emph{grounded} in its requirement,
\emph{actionable} by a test engineer, and its expected result is
\emph{correct}. The judge belongs to the same model family as the generator,
which is a known self-preference risk; the structural matcher was introduced
to remove it from the headline figure.

% ---------------------------------------------------------------------------
\subsection{Method 1: RAG pipeline}
\label{sec:exp-rag}

\paragraph{Method.}
The first method is an agentic retrieval-augmented pipeline organised in
three stages. Axis~A extracts the PDF, splits it with a structure-aware
two-pass chunker and indexes the chunks for hybrid dense and sparse retrieval
fused by reciprocal rank. Axis~B runs a Planner and then a chain of three
agents (Analyst, RAG Router, Designer) once per normative chunk, with
schema-constrained generation, an error-reflective retry and a BERTScore
self-evaluation gate, and assembles the output as an ISO/IEC/IEEE 29119-3
test plan. Axis~C audits each generated case with structural guardrails, an
embedding-distance pre-audit and an Auditor agent whose verdict is overridden
by a deterministic rule. The pipeline was developed with \texttt{llama3.2};
the full run reported here uses \texttt{gemma4:31b} for every agent, served
locally with a context of 4096 tokens.

\paragraph{Results: intermediate artefacts.}
The first documented runs produced chunks and an index but no test plan
(Table~\ref{tab:agentic-baseline}). The measurements of
Table~\ref{tab:diagnosis} show that a single defect, a heading detector that
matched object identifiers in the annexes instead of clause numbers, silently
disabled the chunking, cross-reference, defined-term and planning mechanisms.
A comparison of extraction back-ends (Table~\ref{tab:backends}) located the
cause in the extraction layer.

\paragraph{Results: the full run.}
One full run of the three axes was then made with the heading defect still
present (Table~\ref{tab:agentic-run}). Axis~B wrote a test case for 337 of
the 472 normative chunks (71.4\,\%); the other 135 were dropped because an
agent did not return valid JSON within its retries, the RAG Router alone
accounting for 102 of them. The loss does not depend on the strength of the
requirement (72\,\% of the chunks with a \emph{shall}, 69\,\% with a
\emph{should}, 71\,\% with a \emph{may}) but on the place in the document: it
is highest on the certificate-profile tables. Axis~C passed 154 of the 337
cases (45.7\,\%). Its hallucination and omission rates miss their targets and
its contradiction rate meets it (Table~\ref{tab:agentic-pending}). These rates
are an optimistic estimate: the Auditor is the model that wrote the cases, and
the six cases it failed to audit are counted as free of findings. The schema
check found no violation, and the embedding-distance pre-audit flagged no
case, its score being the same for passed and failed cases (0.84).

The heading defect is visible in the plan. The 1,246 chunks carry 27 distinct
section labels, so the 337 cases share 97 identifiers, 127 of them a single
one, the five feature sets have no meaningful name, and a case can only be
traced to its chunk, not to a clause. For that reason coverage is reported
per normative chunk and not per section, and the run has not been scored
against the official suite: both the judge and the structural matcher need a
clause or a named object on the generated side. The cases are also short,
261 of 337 having exactly two steps.

\begin{table}[t]
  \centering
  \caption{RAG pipeline: full run on ISO/IEC DIS 18013-5
    (\texttt{gemma4:31b}).}
  \label{tab:agentic-run}
  \begin{tabular}{lr}
    \toprule
    Measure & Value \\
    \midrule
    Chunks indexed (Axis A)                 & 1,246 \\
    \quad normative                         & 472 \\
    Distinct section labels                 & 27 \\
    \midrule
    Test cases generated (Axis B)           & 337 / 472 = 71.4\,\% \\
    \quad \emph{shall} / \emph{should} / \emph{may} & 248 / 61 / 28 \\
    Chunks dropped                          & 135 \\
    \quad RAG Router / Designer / Analyst   & 102 / 31 / 2 \\
    Distinct test case identifiers          & 97 \\
    Steps per case, mean                    & 2.3 \\
    \midrule
    Audit verdict (Axis C): pass            & 154 = 45.7\,\% \\
    \quad warning                           & 96 = 28.5\,\% \\
    \quad fail                              & 87 = 25.8\,\% \\
    Cases the Auditor could not audit       & 6 \\
    Structural violations                   & 0 \\
    \midrule
    Time: Axis A / B / C                    & 37\,s / 35.1\,h / 6.3\,h \\
    \bottomrule
  \end{tabular}
\end{table}

\begin{table}[t]
  \centering
  \caption{RAG pipeline: state of the documented runs.}
  \label{tab:agentic-baseline}
  \begin{tabular}{lr}
    \toprule
    Measure & Value \\
    \midrule
    Chunks produced (three successive runs) & 843 / 1,175 / 1,202 \\
    Normative chunks in the indexed run     & 466 \\
    Unit tests (all external calls mocked)  & 23 \\
    Test plans produced by these runs       & 0 \\
    \bottomrule
  \end{tabular}
\end{table}

\begin{table}[t]
  \centering
  \caption{RAG pipeline: measured diagnosis on the 1,202-chunk run.}
  \label{tab:diagnosis}
  \small\setlength{\tabcolsep}{4pt}
  \begin{tabular}{lrl}
    \toprule
    Measurement & Value & Consequence \\
    \midrule
    Distinct section labels                  & 25          & about 200 expected \\
    Chunks labelled with one ASN.1 OID       & 645 (54\,\%) & headings matched on OIDs \\
    Chunks labelled \texttt{8.1.2.2}         & 327 (27\,\%) & no heading for 40 pages \\
    Normative chunks holding a full clause  & 1 / 466     & chunking inactive \\
    Chunks with defined terms                & 0 / 1,202   & term injection inactive \\
    Cross-references resolved                & 22, all to OIDs & resolution inactive \\
    Planner clusters                         & 5           & feature sets meaningless \\
    \bottomrule
  \end{tabular}
\end{table}

\begin{table}[t]
  \centering
  \caption{RAG pipeline: extraction back-ends on ISO/IEC DIS 18013-5.}
  \label{tab:backends}
  \small\setlength{\tabcolsep}{4pt}
  \begin{tabular}{llllr}
    \toprule
    Back-end & Headings recovered & Ligatures & Font data & Time \\
    \midrule
    pdfplumber & 0.89 recall & broken & no  & 60\,s \\
    MarkItDown & none                         & broken & no  & 47\,s \\
    PyMuPDF    & 47 / 47 bookmarks     & intact & yes & 5\,s \\
    \bottomrule
  \end{tabular}
\end{table}

\begin{table}[t]
  \centering
  \caption{RAG pipeline: planned evaluation. Coverage is per normative
    chunk, the section labels being unreliable; the Axis~C rates are over the
    337 generated cases.}
  \label{tab:agentic-pending}
  \small\setlength{\tabcolsep}{4pt}
  \begin{tabular}{llrr}
    \toprule
    Axis & Metric & Target & Result \\
    \midrule
    A & Retrieval precision@5           &             & \tbd \\
    A & Retrieval recall@5              &             & \tbd \\
    A & Retrieval F1@5                  &             & \tbd \\
    \midrule
    B & Test cases generated            &             & 337 \\
    B & Normative chunks covered        & $\geq 80\,\%$ & 71.4\,\% \\
    B & Human score (5 criteria, 1--5)  & $\geq 3.5$    & \tbd \\
    \midrule
    C & Hallucination rate              & $< 10\,\%$    & 23.4\,\% \\
    C & Contradiction rate              & $< 5\,\%$     & 4.7\,\% \\
    C & Omission rate                   & $< 20\,\%$    & 43.3\,\% \\
    C & Adversarial detection rate      & $\geq 80\,\%$ & \tbd \\
    C & Auditor--human agreement (Cohen's $\kappa$) &   & \tbd \\
    \bottomrule
  \end{tabular}
\end{table}

% ---------------------------------------------------------------------------
\subsection{Method 2: deterministic pipeline}
\label{sec:exp-deterministic}

The second method is a fixed pipeline with no embeddings, no retrieval and no
agent framework: the control flow is written in code, the language model is
called only to extract and to write, and every reply is checked by rule.
Whatever fails a check is recorded as an explicit gap instead of being
emitted. Three successive versions of this pipeline were evaluated; they are
reported separately because each changes what the generator enumerates.

\subsubsection{Version 1: sentence-level generation}
\label{sec:exp-sentence}

\paragraph{Method.}
A rule-based ingest segments the document into clauses from its layout;
candidate normative sentences are selected by rule and one LLM call per clause
extracts requirements, each accepted only if its quote is verbatim and its
modality matches; one LLM call per requirement then writes a
nominal case and its negative and boundary variants, checked by a rule-based
verifier with a bounded retry. The model is \texttt{gemma4:31b}.

\paragraph{Results.}
Table~\ref{tab:sentence} gives the scores against the official suite.
Coverage is
4.2\,\%, and the three most-tested clauses score 3\,\%, 1\,\% and 0\,\%. The
generated cases are nevertheless mostly valid, which places the problem in
what is generated rather than in how it is written: the ingest recovered only
117 clauses, and one case per prose sentence does not have the granularity of
the official suite.

\begin{table}[t]
  \centering
  \caption{Deterministic pipeline, version 1 (\texttt{gemma4:31b}).}
  \label{tab:sentence}
  \begin{tabular}{lr}
    \toprule
    Measure & Value \\
    \midrule
    Clauses ingested                 & 117 \\
    Requirements extracted           & 215 \\
    Test cases generated             & 428 \\
    \quad nominal / negative / boundary & 215 / 186 / 27 \\
    \midrule
    Judge coverage, all citing cases & 29 / 687 = 4.2\,\% \\
    \quad Appendix 1                 & 3 / 14 \\
    \quad Appendix 2 (mDL)           & 16 / 451 \\
    \quad Appendix 3 (mDL reader)    & 10 / 222 \\
    \midrule
    Precision: grounded              & 52 / 60 \\
    Precision: actionable            & 50 / 60 \\
    Precision: expected result correct & 51 / 60 \\
    \bottomrule
  \end{tabular}
\end{table}

\subsubsection{Version 2: explicit test model and coverage items}
\label{sec:exp-testmodel}

\paragraph{Method.}
The official suite is an enumeration of objects against a fixed ladder of
checks, so an explicit test model was inserted between extraction and
generation. A second LLM stage extracts the objects the standard defines
(kind, presence, type, value domain, size, relations, each with a verbatim
quote); a deterministic stage derives coverage items from requirements and
objects by equivalence partitioning and boundary analysis; and one LLM call
per coverage item writes the test case, the kind of case being set by the tool
and never by the model. The ingest was extended to keep table rows as records,
and extraction accepts drafts individually instead of rejecting a whole batch.
The model is unchanged.

\paragraph{Results.}
Table~\ref{tab:testmodel} compares this run with version~1 on the same
model. Coverage of the official suite rises from 4.2\,\% to 13.8\,\%, the
best clause reaching 25.4\,\%, and 99.2\,\% of the coverage items produce a
case. Actionability does not move and becomes the quality ceiling. The
remaining limit is measured: twelve of the sixteen object-extraction failures
are truncated replies on the largest table clauses, which hold 361 table rows
and yielded 4 objects.

\begin{table}[t]
  \centering
  \caption{Deterministic pipeline, version 2 against version 1, same model
    (\texttt{gemma4:31b}).}
  \label{tab:testmodel}
  \begin{tabular}{lrr}
    \toprule
    Measure & Version 1 & Version 2 \\
    \midrule
    Clauses ingested            & 117 & 265 \\
    Requirements                & 215 & 842 \\
    Objects in the test model   & --  & 921 \\
    Coverage items              & --  & 3,257 \\
    Test cases                  & 428 & 3,231 \\
    Gaps                        & 14  & 66 \\
    \midrule
    Judge coverage              & 29 / 687 = 4.2\,\% & 95 / 687 = 13.8\,\% \\
    \quad Appendix 1            & 3 / 14   & 3 / 14 \\
    \quad Appendix 2 (mDL)      & 16 / 451 & 69 / 451 \\
    \quad Appendix 3 (mDL reader) & 10 / 222 & 23 / 222 \\
    \midrule
    Precision: grounded         & 52 / 60 & 56 / 60 \\
    Precision: actionable       & 50 / 60 & 50 / 60 \\
    Precision: expected result correct & 51 / 60 & 54 / 60 \\
    \bottomrule
  \end{tabular}
\end{table}

\subsubsection{Version 3: conformance-testing methodology}
\label{sec:exp-conformance}

\paragraph{Method.}
The third version adopts the methodology of the conformance-testing
standards the official suite follows (ISO/IEC 9646-2, ITU-T X.291). Four
elements are added to the pipeline. A deterministic grammar stage harvests
objects from formal declarations (CDDL, ABNF, JSON Schema, ASN.1). An
Implementation Conformance Statement (ICS) proforma is derived, with one item
per declared feature and a status of mandatory, optional or conditional, and
one \emph{capability} coverage item is emitted per mandatory or optional
item. Coverage items are organised in the four test groups of X.291
(capability, valid, invalid and inopportune behaviour), with the invalid side
of the ladder emitted for every object. Generation is batched per source. The
evaluation is rebuilt at the same time, with the judge-free structural matcher
and three explicit denominators. The full run of this version was made with
\texttt{qwen3.8:27b}, not with the model of version~2, so every
model-dependent figure below carries a change of model as well as a change of
method.

\paragraph{Results: the deterministic stages.}
The grammar, ICS and derivation stages use no model, so their output is
exactly reproducible from the input alone (Table~\ref{tab:deterministic}).
About a third of the 94 objects first harvested were malformed; the 152
harvested at the end are all well formed. The
share of due ICS items that carry a coverage item is 100\,\% by construction;
it is a soundness invariant of the derivation, not a score, and the
informative figure is the number of ICS items found.

\begin{table}[t]
  \centering
  \caption{Version 3: deterministic stages alone, no model, on ISO/IEC DIS
    18013-5 (265 clauses).}
  \label{tab:deterministic}
  \small\setlength{\tabcolsep}{4pt}
  \begin{tabular}{lrr}
    \toprule
    Measure & Start of development & End of development \\
    \midrule
    Objects harvested from formal notation & 94 & 152 \\
    ICS items                              & --   & 75 \\
    Coverage items                         & 250  & 604 \\
    \quad after deflation                  & --   & 596 \\
    Invalid-side coverage items            & 0    & 198 \\
    Test groups represented                & 2 of 4 & 4 of 4 \\
    \quad capability / valid / invalid     & --   & 75 / 331 / 190 \\
    Due ICS items with a coverage item     & --   & 75 / 75 \\
    \bottomrule
  \end{tabular}
\end{table}

\paragraph{Results: the full run.}
Table~\ref{tab:conformance-run} compares the full run with version~2.
Everything upstream of generation improved: the table clauses that had
truncated are now read, the object count rises by a factor of 1.73, and the
structural match on the comparable basis rises from 22.5\,\% to 34.9\,\%. The
generation stage, however, returned no test case in this run, so judge
coverage and precision are void, not low, and remain to be measured.

\begin{table}[t]
  \centering
  \caption{Version 3 against version 2. The two runs use different
    models. Groups: capability, valid, invalid and inopportune behaviour.
    G / A / C: grounded, actionable and correct.}
  \label{tab:conformance-run}
  \small\setlength{\tabcolsep}{4pt}
  \begin{tabular}{lrr}
    \toprule
    Measure & Version 2 & Version 3 \\
            & (\texttt{gemma4:31b}) & (\texttt{qwen3.8:27b}) \\
    \midrule
    Requirements                  & 842   & 810 \\
    Objects                       & 921   & 1,590 \\
    \quad of which from grammar   & --    & 152 \\
    ICS items                     & --    & 1,989 \\
    \quad mand. / opt. / cond. & -- & 1,438 / 394 / 157 \\
    Coverage items                & 3,257 & 7,978 \\
    \quad cap. / valid / invalid / inopp. & -- & 1,723 / 4,052 / 2,186 / 17 \\
    Due ICS items with a coverage item & -- & 1,832 / 1,832 \\
    \midrule
    Structural match, objects only & 168 / 747 = 22.5\,\% & 261 / 747 = 34.9\,\% \\
    Structural match, every source & 168 / 747 = 22.5\,\% & 276 / 747 = 36.9\,\% \\
    \midrule
    Test cases                    & 3,231 & 0 \\
    Judge coverage                & 95 / 687 = 13.8\,\% & void \\
    Precision (G / A / C, of 60) & 56 / 50 / 54 & void \\
    \bottomrule
  \end{tabular}
\end{table}

\paragraph{Results: denominators and macro averages.}
The size of the official suite is not its scored size. Of the 907 official
cases, 687 cite a clause of the standard under test, and only 250 cite a
clause identifier that exists in the 2020 draft used as input.
Table~\ref{tab:denominators} reports every rate against the three
denominators, with the macro averages over clauses and over test families
beside the per-case rate. On version~2, where both measures exist, the judge
and the structural matcher agree on 73.4\,\% of the cases both can score
(434 of 591).

\begin{table}[t]
  \centering
  \caption{Coverage of the official suite against three denominators. Micro is
    the per-case rate; macro is the mean of the per-clause or per-family
    rates. In version~2 no requirement carries a name, so the structural
    match is on named objects only.}
  \label{tab:denominators}
  \small\setlength{\tabcolsep}{4pt}
  \begin{tabular}{llrrrr}
    \toprule
    Scope & Run and measure & Scored & Micro & Macro & Macro \\
          &                 &        &       & clauses & families \\
    \midrule
    All official cases & V2, judge                      & 687 & 13.8\,\% & 17.0\,\% & 11.5\,\% \\
    ($n = 907$)        & V2, structural                 & 747 & 22.5\,\% & 20.5\,\% & 17.5\,\% \\
                       & V3, structural, objects only   & 747 & 34.9\,\% & 24.5\,\% & 39.0\,\% \\
                       & V3, structural, every source   & 747 & 36.9\,\% & 28.6\,\% & 40.6\,\% \\
    \midrule
    Citing the standard & V2, judge                     & 687 & 13.8\,\% & 17.0\,\% & 11.5\,\% \\
    ($n = 687$)        & V2, structural                 & 591 & 22.8\,\% & 20.5\,\% & 17.1\,\% \\
                       & V3, structural, objects only   & 591 & 34.7\,\% & 24.2\,\% & 39.9\,\% \\
                       & V3, structural, every source   & 591 & 36.7\,\% & 28.4\,\% & 40.9\,\% \\
    \midrule
    In-edition         & V2, judge                      & 250 & 16.8\,\% & 14.9\,\% & 23.3\,\% \\
    ($n = 250$)        & V2, structural                 & 207 & 18.8\,\% & 10.0\,\% & 14.8\,\% \\
                       & V3, structural, objects only   & 207 & 33.3\,\% & 20.9\,\% & 27.3\,\% \\
                       & V3, structural, every source   & 207 & 35.7\,\% & 32.5\,\% & 33.9\,\% \\
    \bottomrule
  \end{tabular}
\end{table}

\paragraph{Results: by check family and by clause.}
The gain is concentrated where the diagnosis of version~2 placed the limit
(Tables~\ref{tab:families} and~\ref{tab:clauses}): the encoding, consistency
and presence families, and the data-element clauses 7.2 and 7.2.1, which go
together from 40 to 118 matched cases of 166. Behaviour, the largest family of
the suite, stays at 16.4\,\%, size falls to zero, and clauses 9.2.2 and 9.3.2
are not matched at all in either run. In total 145 official cases became
matched and 37 stopped being matched.

\begin{table}[t]
  \centering
  \caption{Structural match by check family, matched of scorable official
    cases.}
  \label{tab:families}
  \begin{tabular}{lrrrr}
    \toprule
    Family & \multicolumn{2}{c}{Version 2} & \multicolumn{2}{c}{Version 3} \\
    \midrule
    Encoding     & 35 / 117 & 29.9\,\% & 85 / 117 & 72.6\,\% \\
    Consistency  & 7 / 31   & 22.6\,\% & 21 / 31  & 67.7\,\% \\
    Presence     & 45 / 137 & 32.8\,\% & 69 / 137 & 50.4\,\% \\
    Value domain & 56 / 165 & 33.9\,\% & 62 / 165 & 37.6\,\% \\
    Behaviour    & 20 / 220 & 9.1\,\%  & 36 / 220 & 16.4\,\% \\
    Reuse        & 2 / 24   & 8.3\,\%  & 3 / 24   & 12.5\,\% \\
    Size         & 3 / 53   & 5.7\,\%  & 0 / 53   & 0.0\,\% \\
    \midrule
    All scorable & 168 / 747 & 22.5\,\% & 276 / 747 & 36.9\,\% \\
    \bottomrule
  \end{tabular}
\end{table}

\begin{table}[t]
  \centering
  \caption{Structural match on the twelve most-tested clauses (2021
    numbering), matched of scorable official cases.}
  \label{tab:clauses}
  \begin{tabular}{lrr}
    \toprule
    Clause & Version 2 & Version 3 \\
    \midrule
    7.2.1       & 16 / 90 & 69 / 90 \\
    7.2         & 24 / 76 & 49 / 76 \\
    9.1.2.5     & 9 / 53  & 17 / 53 \\
    9.1.2.4     & 27 / 82 & 30 / 82 \\
    8.3.2.1.2.2 & 19 / 86 & 22 / 86 \\
    8.2.1.1     & 12 / 49 & 15 / 49 \\
    9.1.4       & 6 / 34  & 9 / 34 \\
    8.3.2.1.2   & 5 / 27  & 6 / 27 \\
    8.3.2.1.2.3 & 9 / 32  & 9 / 32 \\
    9.1.1.4     & 13 / 54 & 8 / 54 \\
    9.2.2       & 0 / 33  & 0 / 33 \\
    9.3.2       & 0 / 30  & 0 / 30 \\
    \bottomrule
  \end{tabular}
\end{table}

% ---------------------------------------------------------------------------
\subsection{Method 3: agentic harness}
\label{sec:exp-harness}

% TODO: complete this subsection once the runs are done.
\paragraph{Method.}
% TODO: describe the protocol. Points to fix: the prompt given to the harness,
% what it may read (the standard only, never TS 18013-6), the output format
% (the same plan schema, so that the same evaluation scripts apply), the
% models, and the judge model, which should be the same for every row.
The third method gives the task to an industrial agentic coding harness,
Claude Code, with no purpose-built pipeline around it. The harness receives
the same input document and is asked to produce a test plan in the same
format, which is then scored with the same evaluation as the deterministic
pipeline. Two families of models are run through the harness: Anthropic
models (\tbd) and open-source models (\tbd).

\paragraph{Results.}
Table~\ref{tab:harness} gives the results per model. \tbd

\begin{table}[t]
  \centering
  \caption{Agentic harness: results per model (pending). Judge: judge
    coverage; Struct.: structural match; G / A / C: grounded, actionable and
    correct, of 60 sampled cases.}
  \label{tab:harness}
  \small\setlength{\tabcolsep}{4pt}
  \begin{tabular}{llrrrrr}
    \toprule
    Harness & Model & Cases & Judge & Struct. & G / A / C & Time \\
    \midrule
    Claude Code & Anthropic, \tbd   & \tbd & \tbd & \tbd & \tbd & \tbd \\
    Claude Code & Anthropic, \tbd   & \tbd & \tbd & \tbd & \tbd & \tbd \\
    Claude Code & Open source, \tbd & \tbd & \tbd & \tbd & \tbd & \tbd \\
    Claude Code & Open source, \tbd & \tbd & \tbd & \tbd & \tbd & \tbd \\
    \bottomrule
  \end{tabular}
\end{table}

% ---------------------------------------------------------------------------
\subsection{Comparison of the three methods}
\label{sec:exp-comparison}

% TODO: write the comparison once the RAG and harness results are in.
% Note: the structural match needs a plan that carries coverage items with a
% named source object; decide how it applies to the RAG and harness outputs.
Table~\ref{tab:comparison} places the three methods side by side on the
metrics of Section~\ref{sec:exp-setting}. \tbd

\begin{table}[t]
  \centering
  \caption{The three methods on ISO/IEC DIS 18013-5 (harness results and
    the scoring of the RAG plan pending). Judge: judge coverage, of the 687 official cases citing the
    standard; Struct.: structural match, of the 747 scorable official cases;
    G / A / C: grounded, actionable and correct, of 60 sampled cases.}
  \label{tab:comparison}
  \small\setlength{\tabcolsep}{4pt}
  \begin{tabular}{llrrrrr}
    \toprule
    Method & Model & Cases & Judge & Struct. & G / A / C & Time \\
    \midrule
    1. RAG pipeline       & \texttt{gemma4:31b}  & 337   & \tbd     & \tbd     & \tbd         & 41.4\,h \\
    \midrule
    2. Deterministic, v1  & \texttt{gemma4:31b}  & 428   & 4.2\,\%  & --       & 52 / 50 / 51 & 0.75\,h \\
    2. Deterministic, v2  & \texttt{gemma4:31b}  & 3,231 & 13.8\,\% & 22.5\,\% & 56 / 50 / 54 & 12.7\,h \\
    2. Deterministic, v3  & \texttt{qwen3.8:27b} & 0     & void     & 36.9\,\% & void         & 126.1\,h \\
    \midrule
    3. Harness            & Anthropic, \tbd      & \tbd  & \tbd     & \tbd     & \tbd         & \tbd \\
    3. Harness            & Open source, \tbd    & \tbd  & \tbd     & \tbd     & \tbd         & \tbd \\
    \bottomrule
  \end{tabular}
\end{table}
 from the main article. Requires \usepackage{booktabs}.
% Every \tbd marks a result that has not been computed yet.
% Sources of the figures: docs/article/outline.md and, under method2-deterministic/,
% docs/evaluation/*.md, eval/eval_remote_1/eval/summary.txt, eval/eval_remote_2/eval/summary.txt.
% Method 1 full run: method1-rag/data/output/vps_run_2026-11-19/ (plan, audit report, run_all.log).
\providecommand{\tbd}{\textit{tbd}}

\section{Experiments and Results}
\label{sec:experiments}

Three methods are compared on the same task, generating a conformance test
plan from the text of a standard: a retrieval-augmented agentic pipeline
(Section~\ref{sec:exp-rag}), a deterministic pipeline in which the language
model only fills the steps a fixed procedure assigns to it
(Section~\ref{sec:exp-deterministic}), and an industrial agentic harness
given the task directly (Section~\ref{sec:exp-harness}). Each subsection
gives the method in a short paragraph and then its results;
Section~\ref{sec:exp-comparison} places the three side by side.

\subsection{Experimental setting}
\label{sec:exp-setting}

\paragraph{Input and ground truth.}
All experiments use the same input, the 126-page draft ISO/IEC DIS
18013-5:2020, and the same ground truth, the three appendices of ISO/IEC TS
18013-6:2025, parsed by rule into structured test cases
(Table~\ref{tab:ground-truth}). The 687 official cases that cite the standard
refer to 54 distinct clauses of the published 2021 edition, whereas the input
is the 2020 draft; this edition mismatch is a confound in every coverage
figure below and is quantified in Section~\ref{sec:exp-conformance}.

\begin{table}[t]
  \centering
  \caption{Ground truth: official test cases of ISO/IEC TS 18013-6:2025.}
  \label{tab:ground-truth}
  \begin{tabular}{lrr}
    \toprule
    Appendix & Test cases & Citing an 18013-5 clause \\
    \midrule
    1 -- Common and certificates & 138 & 14 \\
    2 -- mDL                     & 460 & 451 \\
    3 -- mDL reader              & 309 & 222 \\
    \midrule
    Total                        & 907 & 687 \\
    \bottomrule
  \end{tabular}
\end{table}

\paragraph{Metrics.}
\emph{Judge coverage} is the share of official cases for which an LLM judge,
shown the five generated cases closest to the official one, finds at least one
that verifies the same check (same data item or behaviour, same condition,
same expected outcome). \emph{Structural match} uses no model: each official
case is reduced by rule to a pair (data item, check type), and it is matched
when the generated plan carries a coverage item whose source object is named
by that data item and whose check belongs to the same family; 747 of the 907
official cases can be reduced to such a pair and are called \emph{scorable}.
\emph{Precision} is judged on a fixed random sample of 60 generated cases
against three criteria: the case is \emph{grounded} in its requirement,
\emph{actionable} by a test engineer, and its expected result is
\emph{correct}. The judge belongs to the same model family as the generator,
which is a known self-preference risk; the structural matcher was introduced
to remove it from the headline figure.

% ---------------------------------------------------------------------------
\subsection{Method 1: RAG pipeline}
\label{sec:exp-rag}

\paragraph{Method.}
The first method is an agentic retrieval-augmented pipeline organised in
three stages. Axis~A extracts the PDF, splits it with a structure-aware
two-pass chunker and indexes the chunks for hybrid dense and sparse retrieval
fused by reciprocal rank. Axis~B runs a Planner and then a chain of three
agents (Analyst, RAG Router, Designer) once per normative chunk, with
schema-constrained generation, an error-reflective retry and a BERTScore
self-evaluation gate, and assembles the output as an ISO/IEC/IEEE 29119-3
test plan. Axis~C audits each generated case with structural guardrails, an
embedding-distance pre-audit and an Auditor agent whose verdict is overridden
by a deterministic rule. The pipeline was developed with \texttt{llama3.2};
the full run reported here uses \texttt{gemma4:31b} for every agent, served
locally with a context of 4096 tokens.

\paragraph{Results: intermediate artefacts.}
The first documented runs produced chunks and an index but no test plan
(Table~\ref{tab:agentic-baseline}). The measurements of
Table~\ref{tab:diagnosis} show that a single defect, a heading detector that
matched object identifiers in the annexes instead of clause numbers, silently
disabled the chunking, cross-reference, defined-term and planning mechanisms.
A comparison of extraction back-ends (Table~\ref{tab:backends}) located the
cause in the extraction layer.

\paragraph{Results: the full run.}
One full run of the three axes was then made with the heading defect still
present (Table~\ref{tab:agentic-run}). Axis~B wrote a test case for 337 of
the 472 normative chunks (71.4\,\%); the other 135 were dropped because an
agent did not return valid JSON within its retries, the RAG Router alone
accounting for 102 of them. The loss does not depend on the strength of the
requirement (72\,\% of the chunks with a \emph{shall}, 69\,\% with a
\emph{should}, 71\,\% with a \emph{may}) but on the place in the document: it
is highest on the certificate-profile tables. Axis~C passed 154 of the 337
cases (45.7\,\%). Its hallucination and omission rates miss their targets and
its contradiction rate meets it (Table~\ref{tab:agentic-pending}). These rates
are an optimistic estimate: the Auditor is the model that wrote the cases, and
the six cases it failed to audit are counted as free of findings. The schema
check found no violation, and the embedding-distance pre-audit flagged no
case, its score being the same for passed and failed cases (0.84).

The heading defect is visible in the plan. The 1,246 chunks carry 27 distinct
section labels, so the 337 cases share 97 identifiers, 127 of them a single
one, the five feature sets have no meaningful name, and a case can only be
traced to its chunk, not to a clause. For that reason coverage is reported
per normative chunk and not per section, and the run has not been scored
against the official suite: both the judge and the structural matcher need a
clause or a named object on the generated side. The cases are also short,
261 of 337 having exactly two steps.

\begin{table}[t]
  \centering
  \caption{RAG pipeline: full run on ISO/IEC DIS 18013-5
    (\texttt{gemma4:31b}).}
  \label{tab:agentic-run}
  \begin{tabular}{lr}
    \toprule
    Measure & Value \\
    \midrule
    Chunks indexed (Axis A)                 & 1,246 \\
    \quad normative                         & 472 \\
    Distinct section labels                 & 27 \\
    \midrule
    Test cases generated (Axis B)           & 337 / 472 = 71.4\,\% \\
    \quad \emph{shall} / \emph{should} / \emph{may} & 248 / 61 / 28 \\
    Chunks dropped                          & 135 \\
    \quad RAG Router / Designer / Analyst   & 102 / 31 / 2 \\
    Distinct test case identifiers          & 97 \\
    Steps per case, mean                    & 2.3 \\
    \midrule
    Audit verdict (Axis C): pass            & 154 = 45.7\,\% \\
    \quad warning                           & 96 = 28.5\,\% \\
    \quad fail                              & 87 = 25.8\,\% \\
    Cases the Auditor could not audit       & 6 \\
    Structural violations                   & 0 \\
    \midrule
    Time: Axis A / B / C                    & 37\,s / 35.1\,h / 6.3\,h \\
    \bottomrule
  \end{tabular}
\end{table}

\begin{table}[t]
  \centering
  \caption{RAG pipeline: state of the documented runs.}
  \label{tab:agentic-baseline}
  \begin{tabular}{lr}
    \toprule
    Measure & Value \\
    \midrule
    Chunks produced (three successive runs) & 843 / 1,175 / 1,202 \\
    Normative chunks in the indexed run     & 466 \\
    Unit tests (all external calls mocked)  & 23 \\
    Test plans produced by these runs       & 0 \\
    \bottomrule
  \end{tabular}
\end{table}

\begin{table}[t]
  \centering
  \caption{RAG pipeline: measured diagnosis on the 1,202-chunk run.}
  \label{tab:diagnosis}
  \small\setlength{\tabcolsep}{4pt}
  \begin{tabular}{lrl}
    \toprule
    Measurement & Value & Consequence \\
    \midrule
    Distinct section labels                  & 25          & about 200 expected \\
    Chunks labelled with one ASN.1 OID       & 645 (54\,\%) & headings matched on OIDs \\
    Chunks labelled \texttt{8.1.2.2}         & 327 (27\,\%) & no heading for 40 pages \\
    Normative chunks holding a full clause  & 1 / 466     & chunking inactive \\
    Chunks with defined terms                & 0 / 1,202   & term injection inactive \\
    Cross-references resolved                & 22, all to OIDs & resolution inactive \\
    Planner clusters                         & 5           & feature sets meaningless \\
    \bottomrule
  \end{tabular}
\end{table}

\begin{table}[t]
  \centering
  \caption{RAG pipeline: extraction back-ends on ISO/IEC DIS 18013-5.}
  \label{tab:backends}
  \small\setlength{\tabcolsep}{4pt}
  \begin{tabular}{llllr}
    \toprule
    Back-end & Headings recovered & Ligatures & Font data & Time \\
    \midrule
    pdfplumber & 0.89 recall & broken & no  & 60\,s \\
    MarkItDown & none                         & broken & no  & 47\,s \\
    PyMuPDF    & 47 / 47 bookmarks     & intact & yes & 5\,s \\
    \bottomrule
  \end{tabular}
\end{table}

\begin{table}[t]
  \centering
  \caption{RAG pipeline: planned evaluation. Coverage is per normative
    chunk, the section labels being unreliable; the Axis~C rates are over the
    337 generated cases.}
  \label{tab:agentic-pending}
  \small\setlength{\tabcolsep}{4pt}
  \begin{tabular}{llrr}
    \toprule
    Axis & Metric & Target & Result \\
    \midrule
    A & Retrieval precision@5           &             & \tbd \\
    A & Retrieval recall@5              &             & \tbd \\
    A & Retrieval F1@5                  &             & \tbd \\
    \midrule
    B & Test cases generated            &             & 337 \\
    B & Normative chunks covered        & $\geq 80\,\%$ & 71.4\,\% \\
    B & Human score (5 criteria, 1--5)  & $\geq 3.5$    & \tbd \\
    \midrule
    C & Hallucination rate              & $< 10\,\%$    & 23.4\,\% \\
    C & Contradiction rate              & $< 5\,\%$     & 4.7\,\% \\
    C & Omission rate                   & $< 20\,\%$    & 43.3\,\% \\
    C & Adversarial detection rate      & $\geq 80\,\%$ & \tbd \\
    C & Auditor--human agreement (Cohen's $\kappa$) &   & \tbd \\
    \bottomrule
  \end{tabular}
\end{table}

% ---------------------------------------------------------------------------
\subsection{Method 2: deterministic pipeline}
\label{sec:exp-deterministic}

The second method is a fixed pipeline with no embeddings, no retrieval and no
agent framework: the control flow is written in code, the language model is
called only to extract and to write, and every reply is checked by rule.
Whatever fails a check is recorded as an explicit gap instead of being
emitted. Three successive versions of this pipeline were evaluated; they are
reported separately because each changes what the generator enumerates.

\subsubsection{Version 1: sentence-level generation}
\label{sec:exp-sentence}

\paragraph{Method.}
A rule-based ingest segments the document into clauses from its layout;
candidate normative sentences are selected by rule and one LLM call per clause
extracts requirements, each accepted only if its quote is verbatim and its
modality matches; one LLM call per requirement then writes a
nominal case and its negative and boundary variants, checked by a rule-based
verifier with a bounded retry. The model is \texttt{gemma4:31b}.

\paragraph{Results.}
Table~\ref{tab:sentence} gives the scores against the official suite.
Coverage is
4.2\,\%, and the three most-tested clauses score 3\,\%, 1\,\% and 0\,\%. The
generated cases are nevertheless mostly valid, which places the problem in
what is generated rather than in how it is written: the ingest recovered only
117 clauses, and one case per prose sentence does not have the granularity of
the official suite.

\begin{table}[t]
  \centering
  \caption{Deterministic pipeline, version 1 (\texttt{gemma4:31b}).}
  \label{tab:sentence}
  \begin{tabular}{lr}
    \toprule
    Measure & Value \\
    \midrule
    Clauses ingested                 & 117 \\
    Requirements extracted           & 215 \\
    Test cases generated             & 428 \\
    \quad nominal / negative / boundary & 215 / 186 / 27 \\
    \midrule
    Judge coverage, all citing cases & 29 / 687 = 4.2\,\% \\
    \quad Appendix 1                 & 3 / 14 \\
    \quad Appendix 2 (mDL)           & 16 / 451 \\
    \quad Appendix 3 (mDL reader)    & 10 / 222 \\
    \midrule
    Precision: grounded              & 52 / 60 \\
    Precision: actionable            & 50 / 60 \\
    Precision: expected result correct & 51 / 60 \\
    \bottomrule
  \end{tabular}
\end{table}

\subsubsection{Version 2: explicit test model and coverage items}
\label{sec:exp-testmodel}

\paragraph{Method.}
The official suite is an enumeration of objects against a fixed ladder of
checks, so an explicit test model was inserted between extraction and
generation. A second LLM stage extracts the objects the standard defines
(kind, presence, type, value domain, size, relations, each with a verbatim
quote); a deterministic stage derives coverage items from requirements and
objects by equivalence partitioning and boundary analysis; and one LLM call
per coverage item writes the test case, the kind of case being set by the tool
and never by the model. The ingest was extended to keep table rows as records,
and extraction accepts drafts individually instead of rejecting a whole batch.
The model is unchanged.

\paragraph{Results.}
Table~\ref{tab:testmodel} compares this run with version~1 on the same
model. Coverage of the official suite rises from 4.2\,\% to 13.8\,\%, the
best clause reaching 25.4\,\%, and 99.2\,\% of the coverage items produce a
case. Actionability does not move and becomes the quality ceiling. The
remaining limit is measured: twelve of the sixteen object-extraction failures
are truncated replies on the largest table clauses, which hold 361 table rows
and yielded 4 objects.

\begin{table}[t]
  \centering
  \caption{Deterministic pipeline, version 2 against version 1, same model
    (\texttt{gemma4:31b}).}
  \label{tab:testmodel}
  \begin{tabular}{lrr}
    \toprule
    Measure & Version 1 & Version 2 \\
    \midrule
    Clauses ingested            & 117 & 265 \\
    Requirements                & 215 & 842 \\
    Objects in the test model   & --  & 921 \\
    Coverage items              & --  & 3,257 \\
    Test cases                  & 428 & 3,231 \\
    Gaps                        & 14  & 66 \\
    \midrule
    Judge coverage              & 29 / 687 = 4.2\,\% & 95 / 687 = 13.8\,\% \\
    \quad Appendix 1            & 3 / 14   & 3 / 14 \\
    \quad Appendix 2 (mDL)      & 16 / 451 & 69 / 451 \\
    \quad Appendix 3 (mDL reader) & 10 / 222 & 23 / 222 \\
    \midrule
    Precision: grounded         & 52 / 60 & 56 / 60 \\
    Precision: actionable       & 50 / 60 & 50 / 60 \\
    Precision: expected result correct & 51 / 60 & 54 / 60 \\
    \bottomrule
  \end{tabular}
\end{table}

\subsubsection{Version 3: conformance-testing methodology}
\label{sec:exp-conformance}

\paragraph{Method.}
The third version adopts the methodology of the conformance-testing
standards the official suite follows (ISO/IEC 9646-2, ITU-T X.291). Four
elements are added to the pipeline. A deterministic grammar stage harvests
objects from formal declarations (CDDL, ABNF, JSON Schema, ASN.1). An
Implementation Conformance Statement (ICS) proforma is derived, with one item
per declared feature and a status of mandatory, optional or conditional, and
one \emph{capability} coverage item is emitted per mandatory or optional
item. Coverage items are organised in the four test groups of X.291
(capability, valid, invalid and inopportune behaviour), with the invalid side
of the ladder emitted for every object. Generation is batched per source. The
evaluation is rebuilt at the same time, with the judge-free structural matcher
and three explicit denominators. The full run of this version was made with
\texttt{qwen3.8:27b}, not with the model of version~2, so every
model-dependent figure below carries a change of model as well as a change of
method.

\paragraph{Results: the deterministic stages.}
The grammar, ICS and derivation stages use no model, so their output is
exactly reproducible from the input alone (Table~\ref{tab:deterministic}).
About a third of the 94 objects first harvested were malformed; the 152
harvested at the end are all well formed. The
share of due ICS items that carry a coverage item is 100\,\% by construction;
it is a soundness invariant of the derivation, not a score, and the
informative figure is the number of ICS items found.

\begin{table}[t]
  \centering
  \caption{Version 3: deterministic stages alone, no model, on ISO/IEC DIS
    18013-5 (265 clauses).}
  \label{tab:deterministic}
  \small\setlength{\tabcolsep}{4pt}
  \begin{tabular}{lrr}
    \toprule
    Measure & Start of development & End of development \\
    \midrule
    Objects harvested from formal notation & 94 & 152 \\
    ICS items                              & --   & 75 \\
    Coverage items                         & 250  & 604 \\
    \quad after deflation                  & --   & 596 \\
    Invalid-side coverage items            & 0    & 198 \\
    Test groups represented                & 2 of 4 & 4 of 4 \\
    \quad capability / valid / invalid     & --   & 75 / 331 / 190 \\
    Due ICS items with a coverage item     & --   & 75 / 75 \\
    \bottomrule
  \end{tabular}
\end{table}

\paragraph{Results: the full run.}
Table~\ref{tab:conformance-run} compares the full run with version~2.
Everything upstream of generation improved: the table clauses that had
truncated are now read, the object count rises by a factor of 1.73, and the
structural match on the comparable basis rises from 22.5\,\% to 34.9\,\%. The
generation stage, however, returned no test case in this run, so judge
coverage and precision are void, not low, and remain to be measured.

\begin{table}[t]
  \centering
  \caption{Version 3 against version 2. The two runs use different
    models. Groups: capability, valid, invalid and inopportune behaviour.
    G / A / C: grounded, actionable and correct.}
  \label{tab:conformance-run}
  \small\setlength{\tabcolsep}{4pt}
  \begin{tabular}{lrr}
    \toprule
    Measure & Version 2 & Version 3 \\
            & (\texttt{gemma4:31b}) & (\texttt{qwen3.8:27b}) \\
    \midrule
    Requirements                  & 842   & 810 \\
    Objects                       & 921   & 1,590 \\
    \quad of which from grammar   & --    & 152 \\
    ICS items                     & --    & 1,989 \\
    \quad mand. / opt. / cond. & -- & 1,438 / 394 / 157 \\
    Coverage items                & 3,257 & 7,978 \\
    \quad cap. / valid / invalid / inopp. & -- & 1,723 / 4,052 / 2,186 / 17 \\
    Due ICS items with a coverage item & -- & 1,832 / 1,832 \\
    \midrule
    Structural match, objects only & 168 / 747 = 22.5\,\% & 261 / 747 = 34.9\,\% \\
    Structural match, every source & 168 / 747 = 22.5\,\% & 276 / 747 = 36.9\,\% \\
    \midrule
    Test cases                    & 3,231 & 0 \\
    Judge coverage                & 95 / 687 = 13.8\,\% & void \\
    Precision (G / A / C, of 60) & 56 / 50 / 54 & void \\
    \bottomrule
  \end{tabular}
\end{table}

\paragraph{Results: denominators and macro averages.}
The size of the official suite is not its scored size. Of the 907 official
cases, 687 cite a clause of the standard under test, and only 250 cite a
clause identifier that exists in the 2020 draft used as input.
Table~\ref{tab:denominators} reports every rate against the three
denominators, with the macro averages over clauses and over test families
beside the per-case rate. On version~2, where both measures exist, the judge
and the structural matcher agree on 73.4\,\% of the cases both can score
(434 of 591).

\begin{table}[t]
  \centering
  \caption{Coverage of the official suite against three denominators. Micro is
    the per-case rate; macro is the mean of the per-clause or per-family
    rates. In version~2 no requirement carries a name, so the structural
    match is on named objects only.}
  \label{tab:denominators}
  \small\setlength{\tabcolsep}{4pt}
  \begin{tabular}{llrrrr}
    \toprule
    Scope & Run and measure & Scored & Micro & Macro & Macro \\
          &                 &        &       & clauses & families \\
    \midrule
    All official cases & V2, judge                      & 687 & 13.8\,\% & 17.0\,\% & 11.5\,\% \\
    ($n = 907$)        & V2, structural                 & 747 & 22.5\,\% & 20.5\,\% & 17.5\,\% \\
                       & V3, structural, objects only   & 747 & 34.9\,\% & 24.5\,\% & 39.0\,\% \\
                       & V3, structural, every source   & 747 & 36.9\,\% & 28.6\,\% & 40.6\,\% \\
    \midrule
    Citing the standard & V2, judge                     & 687 & 13.8\,\% & 17.0\,\% & 11.5\,\% \\
    ($n = 687$)        & V2, structural                 & 591 & 22.8\,\% & 20.5\,\% & 17.1\,\% \\
                       & V3, structural, objects only   & 591 & 34.7\,\% & 24.2\,\% & 39.9\,\% \\
                       & V3, structural, every source   & 591 & 36.7\,\% & 28.4\,\% & 40.9\,\% \\
    \midrule
    In-edition         & V2, judge                      & 250 & 16.8\,\% & 14.9\,\% & 23.3\,\% \\
    ($n = 250$)        & V2, structural                 & 207 & 18.8\,\% & 10.0\,\% & 14.8\,\% \\
                       & V3, structural, objects only   & 207 & 33.3\,\% & 20.9\,\% & 27.3\,\% \\
                       & V3, structural, every source   & 207 & 35.7\,\% & 32.5\,\% & 33.9\,\% \\
    \bottomrule
  \end{tabular}
\end{table}

\paragraph{Results: by check family and by clause.}
The gain is concentrated where the diagnosis of version~2 placed the limit
(Tables~\ref{tab:families} and~\ref{tab:clauses}): the encoding, consistency
and presence families, and the data-element clauses 7.2 and 7.2.1, which go
together from 40 to 118 matched cases of 166. Behaviour, the largest family of
the suite, stays at 16.4\,\%, size falls to zero, and clauses 9.2.2 and 9.3.2
are not matched at all in either run. In total 145 official cases became
matched and 37 stopped being matched.

\begin{table}[t]
  \centering
  \caption{Structural match by check family, matched of scorable official
    cases.}
  \label{tab:families}
  \begin{tabular}{lrrrr}
    \toprule
    Family & \multicolumn{2}{c}{Version 2} & \multicolumn{2}{c}{Version 3} \\
    \midrule
    Encoding     & 35 / 117 & 29.9\,\% & 85 / 117 & 72.6\,\% \\
    Consistency  & 7 / 31   & 22.6\,\% & 21 / 31  & 67.7\,\% \\
    Presence     & 45 / 137 & 32.8\,\% & 69 / 137 & 50.4\,\% \\
    Value domain & 56 / 165 & 33.9\,\% & 62 / 165 & 37.6\,\% \\
    Behaviour    & 20 / 220 & 9.1\,\%  & 36 / 220 & 16.4\,\% \\
    Reuse        & 2 / 24   & 8.3\,\%  & 3 / 24   & 12.5\,\% \\
    Size         & 3 / 53   & 5.7\,\%  & 0 / 53   & 0.0\,\% \\
    \midrule
    All scorable & 168 / 747 & 22.5\,\% & 276 / 747 & 36.9\,\% \\
    \bottomrule
  \end{tabular}
\end{table}

\begin{table}[t]
  \centering
  \caption{Structural match on the twelve most-tested clauses (2021
    numbering), matched of scorable official cases.}
  \label{tab:clauses}
  \begin{tabular}{lrr}
    \toprule
    Clause & Version 2 & Version 3 \\
    \midrule
    7.2.1       & 16 / 90 & 69 / 90 \\
    7.2         & 24 / 76 & 49 / 76 \\
    9.1.2.5     & 9 / 53  & 17 / 53 \\
    9.1.2.4     & 27 / 82 & 30 / 82 \\
    8.3.2.1.2.2 & 19 / 86 & 22 / 86 \\
    8.2.1.1     & 12 / 49 & 15 / 49 \\
    9.1.4       & 6 / 34  & 9 / 34 \\
    8.3.2.1.2   & 5 / 27  & 6 / 27 \\
    8.3.2.1.2.3 & 9 / 32  & 9 / 32 \\
    9.1.1.4     & 13 / 54 & 8 / 54 \\
    9.2.2       & 0 / 33  & 0 / 33 \\
    9.3.2       & 0 / 30  & 0 / 30 \\
    \bottomrule
  \end{tabular}
\end{table}

% ---------------------------------------------------------------------------
\subsection{Method 3: agentic harness}
\label{sec:exp-harness}

% TODO: complete this subsection once the runs are done.
\paragraph{Method.}
% TODO: describe the protocol. Points to fix: the prompt given to the harness,
% what it may read (the standard only, never TS 18013-6), the output format
% (the same plan schema, so that the same evaluation scripts apply), the
% models, and the judge model, which should be the same for every row.
The third method gives the task to an industrial agentic coding harness,
Claude Code, with no purpose-built pipeline around it. The harness receives
the same input document and is asked to produce a test plan in the same
format, which is then scored with the same evaluation as the deterministic
pipeline. Two families of models are run through the harness: Anthropic
models (\tbd) and open-source models (\tbd).

\paragraph{Results.}
Table~\ref{tab:harness} gives the results per model. \tbd

\begin{table}[t]
  \centering
  \caption{Agentic harness: results per model (pending). Judge: judge
    coverage; Struct.: structural match; G / A / C: grounded, actionable and
    correct, of 60 sampled cases.}
  \label{tab:harness}
  \small\setlength{\tabcolsep}{4pt}
  \begin{tabular}{llrrrrr}
    \toprule
    Harness & Model & Cases & Judge & Struct. & G / A / C & Time \\
    \midrule
    Claude Code & Anthropic, \tbd   & \tbd & \tbd & \tbd & \tbd & \tbd \\
    Claude Code & Anthropic, \tbd   & \tbd & \tbd & \tbd & \tbd & \tbd \\
    Claude Code & Open source, \tbd & \tbd & \tbd & \tbd & \tbd & \tbd \\
    Claude Code & Open source, \tbd & \tbd & \tbd & \tbd & \tbd & \tbd \\
    \bottomrule
  \end{tabular}
\end{table}

% ---------------------------------------------------------------------------
\subsection{Comparison of the three methods}
\label{sec:exp-comparison}

% TODO: write the comparison once the RAG and harness results are in.
% Note: the structural match needs a plan that carries coverage items with a
% named source object; decide how it applies to the RAG and harness outputs.
Table~\ref{tab:comparison} places the three methods side by side on the
metrics of Section~\ref{sec:exp-setting}. \tbd

\begin{table}[t]
  \centering
  \caption{The three methods on ISO/IEC DIS 18013-5 (harness results and
    the scoring of the RAG plan pending). Judge: judge coverage, of the 687 official cases citing the
    standard; Struct.: structural match, of the 747 scorable official cases;
    G / A / C: grounded, actionable and correct, of 60 sampled cases.}
  \label{tab:comparison}
  \small\setlength{\tabcolsep}{4pt}
  \begin{tabular}{llrrrrr}
    \toprule
    Method & Model & Cases & Judge & Struct. & G / A / C & Time \\
    \midrule
    1. RAG pipeline       & \texttt{gemma4:31b}  & 337   & \tbd     & \tbd     & \tbd         & 41.4\,h \\
    \midrule
    2. Deterministic, v1  & \texttt{gemma4:31b}  & 428   & 4.2\,\%  & --       & 52 / 50 / 51 & 0.75\,h \\
    2. Deterministic, v2  & \texttt{gemma4:31b}  & 3,231 & 13.8\,\% & 22.5\,\% & 56 / 50 / 54 & 12.7\,h \\
    2. Deterministic, v3  & \texttt{qwen3.8:27b} & 0     & void     & 36.9\,\% & void         & 126.1\,h \\
    \midrule
    3. Harness            & Anthropic, \tbd      & \tbd  & \tbd     & \tbd     & \tbd         & \tbd \\
    3. Harness            & Open source, \tbd    & \tbd  & \tbd     & \tbd     & \tbd         & \tbd \\
    \bottomrule
  \end{tabular}
\end{table}
